% referee = single column, 1.75x line spacing for submission
% extra   = loads gji_extra.sty which provides natbib citation commands
\documentclass[extra]{gji}
%\documentclass[referee]{gji}

\usepackage{amsmath}
\usepackage{amssymb}
\usepackage{graphicx}
\graphicspath{{../../}}
\usepackage{booktabs}
\usepackage{pdflscape}
\usepackage[colorlinks=true, allcolors=blue]{hyperref}
% gji_extra.sty embeds natbib but omits \citealt — add it here
\makeatletter
\DeclareRobustCommand\citealt
   {\begingroup\NAT@swafalse\def\NAT@ctype{0}\NAT@parfalse\NAT@citetp}
\makeatother
\usepackage{lineno}
\linenumbers

\DeclareMathOperator{\sgn}{sgn}
\newcommand{\RE}{\mathrm{Re}}
\newcommand{\IM}{\mathrm{Im}}
\newcommand{\infint}{\displaystyle\int_{-\infty}^{+\infty}}

% ====================================================================
\title[Navarrete \& Olive -- RSS inter-rifting model]
  {An analytical recurrent-steady-state model for inter-rifting deformation at magmatic rifts}

\author[Navarrete \& Olive]
  {Iv\'an Navarrete$^1$ \and Jean-Arthur Olive$^1$ \\
  $^1$Laboratoire de G\'eologie, \'Ecole Normale Sup\'erieure,
  PSL Research University, Paris, France
  }

\date{}
\pagerange{\pageref{firstpage}--\pageref{lastpage}}
\pubyear{\the\year{}}

% ====================================================================
\begin{document}
\label{firstpage}
\maketitle

% ====================================================================
\begin{summary}
Inter-rifting deformation, the strain that accumulates between discrete
magmatic episodes, is the least-constrained phase of the rifting cycle. We
model it as the superposed relaxation of a long history of discrete dike
intrusions in an elastic plate overlying a Newtonian viscous half-space.
The non-local coupling of the plate to the half-space turns post-diking
relaxation into a fractional diffusion governed by the half-Laplacian and
by a single characteristic velocity $a = EH/[2\eta(1-\nu^2)]$, where $E$,
$H$ and $\nu$ are the Young's modulus, thickness and Poisson's ratio of the
plate and $\eta$ is the viscosity of the half-space. Deformation spreads
from the rift at the constant speed $a$ with a $1/x$ tail, and Fickian
diffusion is the same relaxation evaluated at a single diffusion time.
Periodic intrusions converge to a recurrent steady state with a closed
form in the Digamma and log-Gamma functions. This state conserves the
opening of each intrusion over one cycle and has a single length scale,
the deformation half-width $aT$, where $T$ is the recurrence interval.
Far-field plate divergence therefore emerges without boundary conditions,
and steady far-field velocities record the memory of past intrusions
rather than continuous creep at depth. A Maxwell visco-elastic half-space
with the elastic modulus of the plate modifies this response only within
about $1.3\,H$ of the axis and during the first few relaxation times.
Applied to post-rifting records from Krafla, Iceland, and Asal, Djibouti,
with $T$ derived from the fitted dike opening and a long-term rift opening
rate, the model favours a viscous over a visco-elastic half-space and
constrains the ratio $EH/\eta$ at both sites. Irregular recurrence,
openings and dike positions spread individual velocity profiles around the
steady state with little bias.
\end{summary}

\begin{keywords}
continental rifting; inter-rifting deformation; fractional diffusion;
visco-elastic relaxation; GPS geodesy; Iceland
\end{keywords}

% ====================================================================
\section{Introduction}
\label{sec:intro}

Continental rifting occurs at divergent plate boundaries, which drive the
transition from stable lithosphere to active oceanic spreading through
rifting cycles that span a broad range of spatial and temporal scales
\citep{corti2009continental,ebinger2013rifting}. These cycles alternate
between discrete rifting events and the inter-rifting periods that
separate them. Each event comprises two phases: elastic strain release
through dike emplacement and normal faulting
\citep{grandin2011seismicity,maccaferri2016time}, followed by post-diking
visco-elastic relaxation that decays on timescales of years to decades
\citep{foulger1992post,nooner2009post}. Secular plate divergence, set by
large-scale mantle dynamics \citep{saria2014present,stamps2021redefining},
persists throughout the cycle: it dominates strain accumulation during
inter-rifting periods and is superposed on the rapid deformation of
rifting events. Geodetic campaigns in Iceland
\citep{heki1993plate,hofton1995horizontal} and the Afar
\citep{vigny2007twenty,nooner2009post,grandin2011seismicity} have
documented the two event phases in considerable detail. The mechanisms
governing the build-up of extensional strain between discrete events,
however, remain the least-constrained aspect of the rifting cycle.

Models of the inter-rifting phase conventionally borrow kinematic
frameworks developed for the interseismic earthquake cycle. The standard
approach adapts the backslip model of \citet{savage1978asthenosphere}: a
normal fault bounding the rift is locked from the surface to a depth $H$
and slips continuously at the plate rate below it. The elastic strain that
accumulates on the locked segment then reproduces the observed near-field
velocity gradient. Alternative formulations represent the diking
rift as a tensile dislocation in an elastic half-space that opens
continuously at a depth-dependent rate \citep{lafemina2005geodetic,konca2018slip},
or prescribe a basal shear traction at the base of the plate to mimic
far-field plate forcing \citep{ali2014geodetic}. These models
share a fundamental physical inconsistency: they ignore the rheological
transition at the brittle--ductile boundary of magmatic rifts, where high
heat flow and partial melt weaken the lower crust and uppermost mantle
\citep{prodehl1997east,plasman2017lithospheric}. Geodetic and seismic
constraints place this transition at 5--15~km depth in Iceland and the
Afar \citep{stevens2021tectonic,weinstein2017fault}. There, the inferred
asthenospheric viscosity ($\eta\approx1$--$5\times10^{18}$~Pa~s) implies
Maxwell relaxation times of a few years \citep{hofton1996post,nooner2009post},
far shorter than the inter-rifting interval and incompatible with the
purely elastic response assumed by backslip models. A continuously opening
dislocation is equally problematic:
dike intrusions are finite, discrete events, and the geodetic record shows
that each episode is followed by a localized transient decaying within a
few years \citep{foulger1992post,reiss2022shallow}.

Models of the post-diking transient itself meet a related obstacle from
two directions. Analytical approaches, exemplified by the diffusion model
of \citet{foulger1992post}, reproduce the broad deformation pattern after
individual intrusions, but they require an a priori viscous layer
thickness, and their solutions return to a constant level far from the
rift (Figure~\ref{fig:hofton_diffusion}). The observed profiles already
capture the transition from post-rifting relaxation to steady
inter-rifting divergence. Twelve to seventeen years after the onset of
rifting, the displacement field traces a broad, arctangent-like shape
rather than a transient confined to the rift, and in the profile of
\citet{hofton1996post} the two ends of the traverse settle onto distinct
levels instead of a common one (Figure~\ref{fig:hofton_diffusion}b).
Numerical visco-elastic simulations
\citep{hofton1995horizontal,hofton1996post,nooner2009post} face an
equivalent difficulty: in any finite domain, velocities must vanish at the
lateral boundaries, which cannot be reconciled with the steady plate
divergence observed in global reference frames
\citep{fernandes2004defining,altamimi2016itrf2014}. The result is a kinematic gap between the
transient post-diking signal and the secular strain rates of global plate
motion models \citep{kreemer2014geodetic,stamps2018virtual}.

\begin{figure*}
  \centering
  \includegraphics[width=\textwidth]{fig/fig_krafla_diffusion_published.png}
  \caption{Post-rifting displacement profiles across the Krafla rift and the
    diffusion models published with them, reproduced without refitting.
    (a) \citet{foulger1992post}, their Figure~2b: displacements of 20
    stations between 1987 and 1990. (b) \citet{hofton1996post}, their
    Figure~5: displacements of 17 stations between 1987 and 1992. Error bars
    are $1\sigma$. Red curves: the best-fit diffusion model of each study,
    $u(x,t) = (b/2)\,\mathrm{erfc}[x/(2\sqrt{\kappa t})]$, written here in
    terms of the full dike opening $b$ (both studies report the dike
    half-width $U_0 = b/2$). Following \citet{foulger1992post}, each
    profile is modelled as the velocity of this solution at a time $t$
    after the intrusions, multiplied by the survey interval, with $t = 11$~yr
    in (a) and $t = 13$~yr in (b), as each study assumes. Published
    parameters: (a) $b = 2.0$~m ($U_0 = 1.0 \pm 0.1$~m) and
    $\kappa = 1.1 \pm 0.3$~m$^2$~s$^{-1}$; (b) $b = 2.54$~m
    ($U_0 = 1.27 \pm 0.14$~m) and $\kappa = 3.11 \pm 0.68$~m$^2$~s$^{-1}$.
    The zero reference of profile (b) is arbitrary \citep{hofton1996post},
    so its curve is shifted by the constant 6.7~cm read from their
    Figure~16, not fitted. Both curves return to a single constant level far
    from the rift, whereas the far-field stations of (b) settle on distinct
    levels on the two flanks, below the curve to the west and above it to
    the east.}
  \label{fig:hofton_diffusion}
\end{figure*}

The theoretical foundations laid for cyclic strike-slip fault deformation
\citep{nur1974postseismic,savage1978asthenosphere,spence1979magma,cohen1984crustal} suggest a different perspective, which we adopt here: inter-rifting
deformation is not a distinct physical process but the emergent
superposition of visco-elastic responses to an infinite history of discrete
dike intrusions. Under this view, the apparent steadiness of inter-rifting velocities arises
naturally from the time-averaging of decaying transients, with no need for
a continuously creeping source.

We develop an analytical model for an elastic crust overlying a viscous
half-space \citep{hetland2006thesis} in which dike intrusions recur periodically at
interval $T$ and each open the rift by $b$, the dike opening. The governing
integro-differential equation reduces, after a Fourier transform, to a
fractional diffusion equation \citep{metzler2000random,samko1993integration} parameterized by a single characteristic velocity
$a = EH/[2\eta(1-\nu^2)]$. In the limit of many repeated intrusions, the
system reaches a Recurrent Steady State (RSS) expressed exactly through the
Digamma and log-Gamma functions, which satisfies the far-field plate
divergence constraint by construction. We then extend the model to a
Maxwell visco-elastic half-space via the elastic--visco-elastic
correspondence principle \citep{hetland2005postseismic,barbot2010unified}, recovering two physically distinct relaxation regimes separated by
the Maxwell timescale $\tau_R = \eta/\mu_m$. We apply both models to two
geodetic datasets: horizontal displacement profiles from the Krafla rift,
Iceland, recorded 12--17 years after the onset of the 1975--1984 rifting
sequence \citep{foulger1992post,hofton1996post}, and a 25-year baseline
time series together with a rift-perpendicular velocity profile from the
Asal rift, Djibouti \citep{vigny2007twenty}. Together, these records
constrain the asthenospheric viscosity of two contrasting settings, show
when a geodetic record can distinguish a viscous from a visco-elastic
half-space, and show that the decadal post-rifting transient and secular
plate divergence are the short- and long-term expressions of a single
process.

% ====================================================================
\section{Elastic-viscous (EV) model}
\label{sec:model}

\subsection{Physical setup and governing equations}
\label{sec:setup}

The model represents the lithosphere as a coupled mechanical system: a thin elastic plate of thickness $H$ (the seismogenic crust) overlying a
Newtonian viscous half-space of viscosity $\eta$ representing the asthenosphere \citep[e.g.,][]{hetland2005postseismic,hetland2006thesis}. The horizontal 
coordinate $x$ is perpendicular to the rift axis; the vertical coordinate $z$ is oriented downward, with the free surface at $z = 0$ 
and the crust--mantle interface at $z = H$. We consider only the horizontal component of surface displacement; the vertical (uplift) field, although
part of the geodetic record \citep{hofton1996post,nooner2009post}, lies outside the scope of this model. Under plane-strain conditions with negligible 
vertical stress ($\sigma_{zz}\approx 0$), the depth-averaged horizontal stress in the elastic layer is related to the horizontal surface displacement 
$u(x,t)$ through Hooke's law:
\begin{equation}
\bar{\sigma}_{xx} = \frac{E}{1-\nu^2}\frac{\partial u}{\partial x},
\label{eq:hookean}
\end{equation}
where $E$ is Young's modulus and $\nu$ is Poisson's ratio. The factor
$(1-\nu^2)^{-1}$ accounts for the lateral confinement of plane strain.
Mechanical equilibrium of the plate requires the horizontal gradient of the
depth-averaged stress to be balanced by the shear traction exerted at its base
by the underlying viscous mantle:
\begin{equation}
H\frac{\partial\bar{\sigma}_{xx}}{\partial x} + \sigma_{xz}^{\text{BOT}} = 0.
\label{eq:equil}
\end{equation}

The rheological character of the asthenosphere enters through the constitutive
relation for the basal traction $\sigma_{xz}^{\text{BOT}}$. For a Newtonian
viscous half-space, the shear stress at $x$ depends on the surface velocity
at \emph{all} other points through the non-local Stokes kernel
\citep{savage1978asthenosphere}:
\begin{equation}
\sigma_{xz}^{\text{BOT}}(x,t) =
  \int_{-\infty}^{+\infty}\frac{2\eta}{\pi(x-x')^2}\,
  \frac{\partial u}{\partial t}(x',t)\,dx',
\label{eq:stokes}
\end{equation}
interpreted as Hadamard finite part. This non-local coupling reflects the
long-range character of viscous flow in the half-space: a velocity perturbation
at $x'$ decays as $(x-x')^{-2}$ rather than being confined to a finite
diffusion radius. In magmatically active rifts, where high heat flow and partial
melt weaken the lower crust \citep{holtzman2016questions,plasman2017lithospheric}, this
coupling operates across the entire plate boundary system.

% \citet{foulger1992post} start from the same force balance but close it
% differently. Rather than coupling the elastic layer to an infinite viscous
% half-space through the non-local Stokes kernel (\ref{eq:stokes}), they place
% a viscous channel of finite thickness $H_v$ over a rigid base. Its basal
% traction is then treated as local, simple-shear drag,
% $\sigma_{xz}^{\mathrm{BOT}} \sim (\eta/H_v)\,\partial u/\partial t$, in place
% of the convolution in (\ref{eq:stokes}). Substituted into (\ref{eq:equil}),
% this closure collapses the governing equation to an ordinary, Fickian
% diffusion equation for the surface displacement, $\partial u/\partial t =
% \kappa\,\partial^2 u/\partial x^2$. The diffusivity $\kappa \propto E H
% H_v/\eta$ this introduces aggregates the elastic-layer properties with the
% assumed viscous-channel thickness $H_v$, rather than constraining a
% half-space viscosity directly. Reproducing a given deformation pattern
% therefore requires an a priori choice of $H_v$ at each site -- the
% limitation already noted in Section~\ref{sec:intro}. Section~\ref{sec:fourier}
% derives what the non-local closure (\ref{eq:stokes}) yields in its place: an
% infinite half-space, rather than a finite channel, replaces the ordinary
% Laplacian of Foulger et al.'s diffusion equation with the half-order
% fractional Laplacian $(-\Delta)^{1/2}$ of Equation~\ref{eq:fractional}. The
% two models therefore differ in the order of the governing diffusion
% equation itself, not only in its diffusivity.

Substituting
(\ref{eq:hookean}) into (\ref{eq:equil}) and using (\ref{eq:stokes}) yields
the governing integro-differential equation for $u$:
\begin{equation}
\frac{EH}{1-\nu^2}\frac{\partial^2 u}{\partial x^2}
+ \int_{-\infty}^{+\infty}\frac{2\eta}{\pi(x-x')^2}
  \frac{\partial u}{\partial t}(x',t)\,dx' = 0.
\label{eq:govPDE}
\end{equation}
Dike intrusions enter as initial conditions: each event opens the rift
instantaneously by $b$, offsetting the block on one side of the axis
relative to the other (Figure~\ref{fig:schematic}).
Section~\ref{sec:step} shows that this offset
subsequently redistributes into a symmetric $\pm b/2$ split about the
pre-intrusion position far from the rift, while the opening measured at the
intrusion site itself remains exactly $b$ at all times.

\begin{figure}
  \centering
  \includegraphics[width=\columnwidth]{fig/EsquemaModelo.png}
  \caption{Model geometry.  The brittle elastic plate (thickness $H$,
    Young's modulus $E$, Poisson's ratio $\nu$) overlies either a Newtonian
    viscous half-space (viscosity $\eta$) or, in its visco-elastic
    extension (Section~\ref{sec:maxwell_model}), a Maxwell half-space of
    the same viscosity $\eta$ that additionally shares the plate's elastic
    moduli $E,\nu$ -- the Elastic-Viscous (EV) and Elastic-Visco-Elastic
    (EVE) cases, respectively.  Dike intrusions open the rift by
    $b$ at intervals $T$; the horizontal surface displacement $u(x,t)$ is
    positive in the $+x$ direction.  The parameter
    $a = EH/[2\eta(1-\nu^2)]$ sets the characteristic velocity that couples
    elastic and viscous deformation.}
  \label{fig:schematic}
\end{figure}

\subsection{Fourier reduction and the characteristic velocity}
\label{sec:fourier}

Equation (\ref{eq:govPDE}) is non-local in $x$ and not tractable in physical
space. A spatial Fourier transform resolves this. Defining the transform pair
\begin{equation*}
\hat{u}(k,t) = \int_{-\infty}^{+\infty}u(x,t)e^{-ikx}\,dx,
\end{equation*}
\begin{equation*}
u(x,t) = \frac{1}{2\pi}\int_{-\infty}^{+\infty}\hat{u}(k,t)e^{ikx}\,dk,
\end{equation*}
the transform of $\partial^2 u/\partial x^2$ is $-k^2\hat{u}$. The transform
of the Stokes kernel convolution is:
\begin{equation}
\mathcal{F}\!\left\{
  \int_{-\infty}^{+\infty}\frac{2\eta}{\pi(x-x')^2}
  \frac{\partial u}{\partial t}\,dx'
\right\} = -2\eta|k|\,\frac{\partial\hat{u}}{\partial t}.
\label{eq:stokes_FT}
\end{equation}
Applying the Fourier transform to (\ref{eq:govPDE}) and using
(\ref{eq:stokes_FT}):
\begin{equation}
-\frac{EH}{1-\nu^2}k^2\hat{u} - 2\eta|k|\frac{\partial\hat{u}}{\partial t} = 0.
\label{eq:fourier_pre}
\end{equation}
Rearranging and introducing the characteristic velocity
\begin{equation}
a := \frac{EH}{2\eta(1-\nu^2)},
\label{eq:a}
\end{equation}
equation (\ref{eq:fourier_pre}) simplifies to a decoupled first-order ODE for
each wavenumber $k$:
\begin{equation}
\frac{\partial\hat{u}}{\partial t} = -a\,|k|\,\hat{u}.
\label{eq:fourier}
\end{equation}
The parameter $a$ encodes the full elastic--viscous coupling of the system:
large $a$ (stiff crust, low viscosity) implies rapid redistribution of
post-diking strain over broad spatial scales; small $a$ implies slow, locally
confined relaxation. In physical space, the operator $|k|$ corresponds to the
fractional half-Laplacian $(-\Delta)^{1/2}$, so (\ref{eq:fourier}) is
equivalently written as the fractional diffusion equation
\citep{metzler2000random,samko1993integration,samko1993fractional}
\begin{equation}
\frac{\partial u}{\partial t} = -a\,(-\Delta)^{1/2}u.
\label{eq:fractional}
\end{equation}
% This order is not a modelling choice. The fractional Laplacian admits a
% representation as a Dirichlet-to-Neumann map for a harmonic extension
% problem into an auxiliary half-space, weighted by $y^{1-2s}$; the weight
% vanishes, and the extension becomes the ordinary Laplace equation, only at
% $s=1/2$ \citep{caffarelli2007extension,daoud2022fractional}. That is the
% same half-space potential problem generating the Stokes kernel itself, so
% the half-order non-locality in (\ref{eq:fractional}) is forced by the
% geometry of the elastic-plate--viscous-half-space coupling, not fitted to
% the data. Solving the full two-dimensional Stokes flow beneath the plate,
% governed by the biharmonic rather than the harmonic equation, still
% recovers the same interface operator $-2\eta(-\Delta)^{1/2}$ exactly once
% the interface stress is extracted.

% This contrasts with classical Fickian diffusion: the half-order operator, a
% direct consequence of the non-local Stokes kernel, produces Lorentzian rather
% than Gaussian spatial profiles, with tails that decay as $1/x$ rather than
% exponentially.

The solution of (\ref{eq:fourier}) for arbitrary initial data is:
\begin{equation}
\hat{u}(k,t) = \hat{u}(k,0)\,e^{-a|k|t}.
\label{eq:fourier_sol}
\end{equation}

\subsection{Step response to a single dike intrusion}
\label{sec:step}

The simplest building block of the model is the deformation produced by a
single instantaneous dike intrusion at $t = 0$, which opens the rift by $b$.
Because the governing equation (\ref{eq:govPDE}) constrains $u$ only through
its gradient, the reference level from which this offset is measured is
arbitrary. We therefore write the initial condition, without loss of
generality, as a Heaviside displacement, with the symmetric convention
$\mathcal{H}(0) := 1/2$ so that the profile is well defined at the rift
axis:
\begin{equation}
u(x,0^+) = b\,\mathcal{H}(-x) = \begin{cases}
  0 \; , & x>0 \\
  b/2 \; , & x=0 \\
  b \; , & x<0
\end{cases}.
\label{eq:IC}
\end{equation}
Its Fourier transform, evaluated in the distributional sense in
Appendix~\ref{app:step}, is
\begin{equation}
\hat{u}(k,0) = b\pi\delta(k) + \frac{ib}{k}.
\label{eq:IC_FT}
\end{equation}
Substituting (\ref{eq:IC_FT}) into (\ref{eq:fourier_sol}) and noting that
$\delta(k)$ is active only where $e^{-a|k|t}=1$:
\begin{equation}
\hat{u}(k,t) = b\pi\delta(k) + \frac{ib}{k}\,e^{-a|k|t}.
\label{eq:evolved}
\end{equation}
The inverse transform of (\ref{eq:evolved}) reduces to a Dirichlet-type
integral (Appendix~\ref{app:step}). We label quantities of this
single-event solution with a STEP superscript, to distinguish them from the
periodic Recurrent Steady State (RSS) of \S\ref{sec:css}. The step
displacement is
\begin{equation}
u^{\text{STEP}}(x,t) = \frac{b}{\pi}\arctan\!\left(\frac{at}{x}\right).
\label{eq:step}
\end{equation}
Differentiating with respect to $t$, the step velocity is a Lorentzian
profile that broadens in time:
\begin{equation}
v^{\text{STEP}}(x,t) = \frac{b a}{\pi}\,\frac{x}{x^2 + a^2t^2}.
\label{eq:stepv}
\end{equation}
Both expressions depend on $x$ and $t$ only through a single combination,
$x/(at)$: the self-similarity structure of the step response, developed
in \S\ref{sec:step_selfsimilar} below.

The surface profile is only the boundary trace of a flow that fills the
underlying half-space. Depth is reintroduced through the same harmonic
extension used to construct the fractional Laplacian in
Section~\ref{sec:fourier}: at depth $z$, the horizontal and vertical
interior velocity components are the Poisson-kernel convolution of the
surface velocity $V(x)$,
\begin{equation}
\begin{split}
v_x(x,z) &= \frac{2}{\pi}\int_{-\pi/2}^{\pi/2} \sin^2\theta\, V(x-z\tan\theta)\,d\theta, \\
v_z(x,z) &= \frac{2}{\pi}\int_{-\pi/2}^{\pi/2} \sin\theta\cos\theta\, V(x-z\tan\theta)\,d\theta,
\end{split}
\label{eq:halfspace_kernel}
\end{equation}
with $V$ set to the EV step velocity $v^{\text{STEP}}(x,t)$ (\ref{eq:stepv}).

Figure~\ref{fig:viscous_halfspace_step} follows the relaxation of a
single intrusion over three orders of magnitude in time, at the surface
(left column) and in the underlying half-space (middle and right
columns). At the surface, the response is a pulse of divergent velocity
that travels away from the rift axis. Its peak sits exactly at
$x = a\tau$ (dashed lines), so $a$ is the speed at which post-diking
deformation propagates, and its amplitude decays as $b/(2\pi\tau)$.
A point at distance $|x|$ has moved by $b/4$, half of its final share
of the opening, when the peak passes it at $\tau = |x|/a$; the remaining
$b/4$ accumulates only slowly, as $u^{\text{STEP}}$ approaches $b/2$
as $1/\tau$. Ahead of the peak ($|x| \gg a\tau$), Equation~(\ref{eq:stepv})
reduces to the time-independent tail $ba/(\pi x)$: a distant station
records a nearly constant velocity from the moment of intrusion until
the peak arrives. Behind the peak ($|x| \ll a\tau$), the velocity falls
as $bx/(\pi a\tau^2)$ and vanishes on the axis itself.

The interior field shows the flow that drives this surface pulse.
Material rises beneath the axis and turns outward beneath the plate,
feeding the divergence of the two flanks (middle column). Because
Equation~(\ref{eq:halfspace_kernel}) is linear
in the surface velocity, the interior field inherits the self-similarity
of Equation~(\ref{eq:stepv}): $\tau$ times the velocity depends only on
$x/(a\tau)$ and $z/(a\tau)$. The deforming region therefore deepens at
the same rate $a$ as it widens. At distances much larger than $a\tau$
from the intrusion, the flow keeps its early, time-independent pattern,
which is why the middle column barely changes across the six rows: even
at $\tau = 3\tau_R$, $a\tau \approx 4H$ remains small compared with the
$30H$ window. Within $a\tau$, the flow has already slowed down. The
enlarged panels (right column) show this as a low-velocity region that
opens beneath the axis and grows with $\tau$. The surface itself moves
only horizontally ($v_z = 0$ at $z = 0$), consistent with the
horizontal-only plate of Section~\ref{sec:setup}; the uplift recorded
geodetically above rifts lies outside this model.
\begin{figure*}
  \centering
  \includegraphics[width=\textwidth]{fig/fig_halfspace_step_contour.png}
  \caption{EV STEP response to a single dike intrusion of opening $b$ at
    $t=0$, at six times $\tau/\tau_R \in \{0.001,\,0.25,\,0.5,\,1,\,2,\,3\}$,
    one per row. The Maxwell time $\tau_R=\eta/\mu_m$
    (Section~\ref{sec:maxwell_model}) serves here only as a time unit, the
    same one used for the EVE counterpart
    (Figure~\ref{fig:maxwell_halfspace_step}).
    Left column: surface velocity $v^{\text{STEP}}$
    (Equation~\ref{eq:stepv}; solid, left axis, in units of $b/\tau_R$)
    and displacement $u^{\text{STEP}}$ (Equation~\ref{eq:step}; dashed,
    right axis, in units of $b$), counted from the instant after the
    intrusion, so that the two flanks move apart symmetrically and each
    tends to $\pm b/2$ at the axis. The dashed vertical line marks
    $x=a\tau$, the position of the velocity peak.
    Middle column: speed $|v|$ in the half-space beneath the plate (filled
    contours, logarithmic scale, in units of $b/\tau_R$), computed from
    the surface velocity with Equation~(\ref{eq:halfspace_kernel}) over
    $|x|\le 30H$ and $z\le 30H$. Arrows show the flow direction only and
    have a fixed length.
    Right column: the same field enlarged to $|x|\le H$, $z\le H$ beneath
    the rift axis, with the same colour scale. Speeds outside the colour
    bar range $[0.01,\,1]$ are clipped to its end colours.
    All lengths are scaled by the plate thickness $H$. Parameters:
    $H=10$~km, $E=40$~GPa, $\nu=0.25$, $\eta=5.2\times10^{18}$~Pa\,s,
    giving $a\approx1.3$~km\,yr$^{-1}$ and $a\tau_R/H=1/(1-\nu)\approx1.33$.}
  \label{fig:viscous_halfspace_step}
\end{figure*}

\subsubsection{Self-similarity of the step response}
\label{sec:step_selfsimilar}

Both the displacement (\ref{eq:step}) and velocity (\ref{eq:stepv}) of
Section~\ref{sec:step} depend on $x$ and $t$ only through the
self-similarity variable $\xi := x/(at)$, which measures distance in units
of the velocity-peak position $at$, so that the profile's shape is invariant
in $\xi$ at every $t$. Dividing $u^{\text{STEP}}$ by $b$ and
$v^{\text{STEP}}$ by $b/t$ collapses both onto universal curves,
independent of $E$, $H$, $\eta$, or $b$:
\begin{equation}
\frac{u^{\text{STEP}}}{b} = \frac{1}{\pi}\arctan\!\left(\frac{1}{\xi}\right),
\label{eq:step_dimless}
\end{equation}
\begin{equation}
\frac{t\,v^{\text{STEP}}}{b} = \frac{1}{\pi}\,\frac{\xi}{1+\xi^2}.
\label{eq:stepv_dimless}
\end{equation}

Figure~\ref{fig:step_selfsimilar}a plots the displacement collapse
(\ref{eq:step_dimless}) directly, alongside the Fickian
\citet{foulger1992post} diffusion mode: writing $u_{\text{Foulger}}/b = \tfrac{1}{2}\,\mathrm{erfc}(\xi_F/2)$ in
terms of its own self-similarity variable $\xi_F := x/\sqrt{\kappa t}$,
both curves are shown for $\xi>0$ only. This isolates the decay of the
tails, a power law for EV and a Gaussian for Foulger, without reference
to any particular parameter.

The same collapse applies to the velocity (\ref{eq:stepv_dimless}).
Differentiating $u_{\text{Foulger}}/b$ in $t$ gives the Fickian model's
velocity collapse:
\begin{equation}
\frac{t\,v_{\text{Foulger}}}{b} = \frac{\xi_F}{4\sqrt{\pi}}\,
  e^{-\xi_F^2/4}.
\label{eq:foulger_v_dimless}
\end{equation}
Figure~\ref{fig:step_selfsimilar}b plots (\ref{eq:stepv_dimless}) and
(\ref{eq:foulger_v_dimless}) together: the EV velocity decays as the power law $\xi^{-1}$,
a straight line on logarithmic axes, while the Fickian velocity decays as
the far steeper $\xi_F\,e^{-\xi_F^2/4}$. Because $\xi$ and $\xi_F$ scale
distance by different lengths, $at$ and $\sqrt{\kappa t}$, the horizontal
offset between the two curves carries no physical meaning; the robust
contrast is the shape of the tail, a power law against a Gaussian.

The two collapses also differ in how the velocity peak moves. The EV peak
sits at $\xi = 1$, that is at $x = at$, and migrates outward at the
constant speed $a$. The Fickian peak sits at $\xi_F = \sqrt{2}$, that is at
$x = \sqrt{2\kappa t}$, and decelerates as $t^{-1/2}$. A station at
distance $x$ from the axis therefore records its maximum velocity at
$t = x/a$ in the EV model but at $t = x^2/(6\kappa)$ in the Fickian
model: the arrival time of the transient grows linearly with distance in
the first case and quadratically in the second. Geodetic time series
recorded at several distances from the axis can, in principle,
discriminate between the two on this basis alone, independently of the
tail amplitude.
\begin{figure}
  \centering
  \includegraphics[width=\columnwidth]{fig/fig_step_selfsimilar_raw.png}
  \caption{Self-similarity collapse of the step response in the EV model
    and in the Fickian \citet{foulger1992post} model, each plotted against
    its own self-similarity variable ($\xi = x/(at)$ for EV,
    $\xi_F = x/\sqrt{\kappa t}$ for Foulger) for $\xi>0$. Both models start
    from the same opening $b$, so both displacements equal $b/2$ at the
    rift axis; the curves are not normalised.
    (a)~Displacement $u/b$: EV (\ref{eq:step_dimless}) and Foulger
    $\tfrac{1}{2}\mathrm{erfc}(\xi_F/2)$.
    (b)~Velocity $t\,v/b$: EV (\ref{eq:stepv_dimless}) and Foulger
    (\ref{eq:foulger_v_dimless}).
    In both panels the dotted guide marks the $\xi^{-1}$ power-law slope
    the EV curve approaches; the Fickian curves instead decay as a
    Gaussian. Because $\xi$ and $\xi_F$ use different length scales, only
    the shapes of the two tails, not their horizontal offset, are
    comparable.}
  \label{fig:step_selfsimilar}
\end{figure}

\subsection{Cyclic accumulation and the Recurrent Steady State}
\label{sec:css}

\subsubsection{Finite sum over past intrusions}

Consider a rift that has hosted $N$ dike intrusions at regular intervals
$T$, the most recent a time $\tau\in[0,T)$ ago. Indexing the intrusions by
$j=0,1,\dots,N-1$, the number of cycles \emph{before} the most recent one,
the $j$-th intrusion has age $\tau+jT$. Superposing the step responses
(\ref{eq:step}), each counted from the instant after its own intrusion,
gives the accumulated displacement:
\begin{equation}
u(x,\tau) = \frac{b}{\pi}\sum_{j=0}^{N-1}
  \arctan\!\left(\frac{a(\tau+jT)}{x}\right).
\label{eq:finite_sum}
\end{equation}
Differentiating with respect to $t$, the corresponding velocity field is:
\begin{equation}
v(x,\tau) = \frac{b a}{\pi}\sum_{j=0}^{N-1}
  \frac{x}{x^2 + a^2(\tau+jT)^2}.
\label{eq:finite_sum_v}
\end{equation}

\subsubsection{The RSS limit and Digamma representation}

After many cycles the system enters a periodic state, the rift analog of the
cyclic steady state described for the interseismic earthquake cycle \citep{thatcher1984viscoelastic,hetland2006interseismic}.
The Recurrent Steady State (RSS) velocity is obtained by extending
(\ref{eq:finite_sum_v}) to an infinite history by letting $N\to\infty$:
\begin{equation}
v_{\text{RSS}}(x,\tau) = \frac{b a}{\pi}
  \sum_{j=0}^{\infty}\frac{x}{x^2 + a^2(\tau+jT)^2}.
\label{eq:css_sum}
\end{equation}
Because this series converges absolutely, it reduces to a closed form in
the Digamma function $\psi(z)$ (Appendix~\ref{app:digamma}):
\begin{equation}
v_{\text{RSS}}(x,\tau) =
  \frac{b}{\pi T}\,\IM\!\left\{
    \psi\!\left(\frac{\tau}{T}+i\frac{x}{aT}\right)
  \right\}.
\label{eq:vcss}
\end{equation}

\subsubsection{RSS displacement and far-field conservation}

Integrating (\ref{eq:vcss}) over $\tau$ from the most recent intrusion,
and using the antiderivative $d[\ln\Gamma(z)]/dz = \psi(z)$, the RSS
displacement accumulated since that intrusion is:
\begin{equation}
u_{\text{RSS}}(x,\tau) =
  \frac{b}{\pi}\,\IM\!\left\{
    \ln\Gamma\!\left(\frac{\tau}{T}+i\frac{x}{aT}\right)
    -\ln\Gamma\!\left(i\frac{x}{aT}\right)
  \right\}.
\label{eq:ucss}
\end{equation}
Over one complete cycle, the cumulative displacement at any fixed $x>0$ equals
half the dike opening:
\begin{equation}
\int_0^T v_{\text{RSS}}(x,\tau)\,d\tau = \frac{b}{2}, \qquad \forall\,x>0.
\label{eq:conservation}
\end{equation}
This is not imposed as a boundary condition but follows analytically from the
Digamma representation: $\Gamma(1+iq) = iq\,\Gamma(iq)$ implies
$\IM\{\ln\Gamma(1+iq)-\ln\Gamma(iq)\} = \pi/2$, giving the integral $b/2$.
The two flanks together therefore separate by $b$ per cycle, exactly the
opening of the dike intruded during that cycle, while each flank advances
by only half of it. The RSS thus reconciles high-frequency magmatic
pulses with secular divergence without unphysical boundary conditions,
satisfying the far-field constraint of global plate motion models
\citep{saria2014present,stamps2021redefining} by construction.

Figure~\ref{fig:rss_accumulation} shows how this conservation plays out
over successive cycles.  After $n$ complete cycles and a further phase
$\tau$, the displacement accumulated since a reference intrusion is
$n(b/2)\,\mathrm{sgn}(x) + u_{\text{RSS}}(x,\tau)$.  At the end of each
cycle, $u_{\text{RSS}}(x,T^-) = b/2$ at every $x>0$
(Equation~\ref{eq:conservation}), so each flank has translated as a rigid
block by $b/2$ and the profile is uniform (thick lines in
Figure~\ref{fig:rss_accumulation}a).  Within the cycle, the displacement is
not uniform: the region next to the axis takes up its share of the
opening almost immediately after the intrusion, while distant points
catch up later in the cycle.  The resulting histories
(Figure~\ref{fig:rss_accumulation}b) range from a staircase at the axis,
which jumps by $b/2$ at each intrusion, to a nearly uniform rate $b/(2T)$
far from it.  At $x = 0.1\,aT$, nearly 80 per cent of the $b/2$ is gained
in the first quarter of the cycle; at $x = aT$, the history departs from
the uniform rate by at most $0.04\,b$, and at $x = 3\,aT$ by at most
$0.013\,b$.  A station at $|x| \gtrsim aT$ therefore records an almost
steady velocity, the inter-rifting plate divergence, even though all of
the opening is delivered in discrete events.

\begin{figure}
  \centering
  \includegraphics[width=\columnwidth]{fig/fig_rss_accumulation.png}
  \caption{Displacement accumulated in the EV RSS over three cycles,
    starting from an intrusion at $t=0$ that follows an infinite periodic
    history.  Displacement is normalised by $b$, distance by $aT$ and time
    by $T$.  (a)~Profiles every $T/4$; thick lines mark the end of each
    cycle, when both flanks have moved uniformly by $b/2$ per cycle.
    (b)~Displacement histories at fixed distances from the axis ($x>0$),
    from the staircase at $x=0^+$ to the near-uniform rate at $x=3\,aT$;
    the dotted line marks the mean rate $(b/2)\,t/T$.}
  \label{fig:rss_accumulation}
\end{figure}

The interior flow of the RSS follows from the same Poisson-kernel
extension (\ref{eq:halfspace_kernel}), with $V$ now set to the EV RSS
velocity $v_{\text{RSS}}(x,\tau)$ (\ref{eq:vcss}).
Figure~\ref{fig:viscous_halfspace_css} shows it in the same layout as the
STEP field of Section~\ref{sec:step} (Figure~\ref{fig:viscous_halfspace_step}),
for $T = 5\tau_R$ and at six phases spread over one cycle: the surface
profile (left column), the speed in the half-space (middle column), and
an enlargement beneath the axis (right column).  Distances are scaled by
$aT$ and velocities by the rift opening rate $V = b/T$, so colours are not
comparable with those of the STEP figure, whose unit is $b/\tau_R$.

The RSS field differs from the STEP field in one essential respect: it
never empties out. Far from the axis ($|x| \gg aT$), each flank moves at
$V/2$ throughout the cycle, so the two flanks separate at the full rift
opening rate $V$, as required by Equation~(\ref{eq:conservation}). Near the axis, each
cycle repeats the behaviour of the single event. Immediately after an
intrusion, the near-axis velocity exceeds $V$ by more than an order of
magnitude, and a velocity peak migrates away from the axis close to
$x = a\tau$ (dashed lines), as in Section~\ref{sec:step}. By the end of
the cycle, the near-axis velocity has fallen well below $V/2$, and the
profile has smoothed into a monotonic ramp from $-V/2$ to $+V/2$ over a
distance of order $aT$. The interior field follows the same pattern.
Upwelling beneath the axis and outward flow beneath the plate persist at
every phase, and over most of the half-space the speed stays between
roughly $0.1\,V$ and $V/2$ throughout the cycle (middle column).  Only the
region beneath the axis changes: as in the STEP response, a low-velocity
region opens there after each intrusion and grows until the next one
(right column).

\begin{figure*}
  \centering
  \includegraphics[width=\textwidth]{fig/fig_halfspace_css_contour.png}
  \caption{EV RSS for $T=5\tau_R$, in the same layout as the STEP response
    (Figure~\ref{fig:viscous_halfspace_step}), at six phases
    $\tau/T \in \{0.001,\,0.2,\,0.4,\,0.6,\,0.8,\,0.99\}$, one per row.
    Left column: surface velocity $v_{\text{RSS}}$
    (Equation~\ref{eq:vcss}; solid, left axis, in units of $V=b/T$) and
    displacement accumulated since the most recent intrusion
    $u_{\text{RSS}}$ (Equation~\ref{eq:ucss}; dashed, right axis, in units
    of $b$).  The dashed vertical line marks $x=a\tau$, the velocity-peak
    position of the most recent intrusion alone.
    Middle column: speed $|v|$ in the half-space (filled contours,
    logarithmic scale, in units of $V$), computed with
    Equation~(\ref{eq:halfspace_kernel}) over $|x|\le 4.5\,aT$ and
    $z\le 4.5\,aT$.  Arrows show the flow direction only and have a fixed
    length.
    Right column: the same field enlarged to $|x|\le 0.15\,aT$,
    $z\le 0.15\,aT$ beneath the rift axis, with the same colour scale.
    Speeds outside the colour bar range $[0.01,\,1]$ are clipped to its
    end colours.  Parameters as in Figure~\ref{fig:viscous_halfspace_step},
    giving $aT \approx 67$~km.}
  \label{fig:viscous_halfspace_css}
\end{figure*}

Equation~(\ref{eq:vcss}) also fixes the scales of the RSS. Normalised by
the rift opening rate $V=b/T$, the velocity depends only on the phase $\tau/T$
and on $x/(aT)$, which together form the Digamma argument
$\tau/T + ix/(aT)$: $aT$ sets the spatial scale
of the deformation zone and $b/T$ its amplitude. Late in the cycle, $aT$
is the distance over which the velocity rises from zero at the axis to
the far-field value $\pm V/2$; early in the cycle, the near-axis profile
is dominated by the most recent intrusion, whose velocity peak sits at
$x \approx a\tau < aT$ (Figure~\ref{fig:viscous_halfspace_css}). Because
$aT \propto EHT/\eta$, a geodetic profile that resolves this scale
constrains $EH/\eta$ wherever $T$ can be estimated independently.

The RSS assumes an infinite history of intrusions. Truncating the sum
(\ref{eq:finite_sum_v}) after $N$ intrusions removes the contribution of
the older ones; approximating the missing terms by an integral, the
truncated velocity falls short of the RSS by
$\approx (V/\pi)\arctan[|x|/(NaT)]$. The truncated sum is therefore
accurate only within $|x| \ll NaT$ (Figure~\ref{fig:css_behavior}): after
ten cycles, the deficit is 3 per cent of $V$ at $|x| = aT$ but 25 per cent
at $|x| = 10\,aT$. Physically, the far-field plateau $\pm V/2$ is carried
by old intrusions, whose single-event velocity peaks have migrated to
$x \approx a(\tau+jT)$; a station at distance $|x|$ records the RSS only
if the rift has hosted intrusions for much longer than $|x|/a$. Where this
condition fails, as on a newly initiated rift segment, the explicit sum
(\ref{eq:finite_sum}) should replace the RSS.

\begin{figure}
  \centering
  \includegraphics[width=\columnwidth]{fig/fig_css_behavior.png}
  \caption{Convergence of the finite-sum velocity
    (\ref{eq:finite_sum_v}) to the RSS (\ref{eq:vcss}) at phase
    $\tau/T=0.5$, for $N = 1$, 3, 10 and 2000 past intrusions (solid curves)
    and the RSS (dashed). Velocity is normalised by $b/T$ and distance by
    $aT$. The deficit of the truncated sum grows with distance,
    approximately as $(b/\pi T)\arctan[x/(NaT)]$: with $N=10$, the sum
    matches the RSS near the axis but remains below the far-field plateau
    $b/(2T)$ (dotted) for $x \gtrsim aT$.}
  \label{fig:css_behavior}
\end{figure}



\section{Visco-elastic (EVE) extension}
\label{sec:eve_model}

We refer to the two half-space rheologies compared throughout this paper as
Elastic-Viscous (EV), for the Newtonian viscous half-space of
Section~\ref{sec:model}, and Elastic-Visco-Elastic (EVE), for the Maxwell
visco-elastic half-space developed in this section.

\subsection{Maxwell half-space and the visco-elastic Stokes kernel}
\label{sec:maxwell_model}

The Maxwell relaxation time $\tau_R := \eta/\mu_m$, where $\mu_m$ is the
mantle shear modulus, separates two regimes: on timescales $t \ll \tau_R$
the asthenosphere responds elastically, and on timescales $t \gg \tau_R$ it
flows viscously \citep{burgmann2008rheology}. We model this by replacing
the Newtonian half-space with a Maxwell body, a Hookean spring and a
Newtonian dashpot in series. By the elastic--visco-elastic correspondence
principle \citep{hetland2005postseismic,barbot2010unified}, a Maxwell body
responds in the Laplace domain like a dashpot with the frequency-dependent
effective viscosity (Appendix~\ref{app:correspondence})
\begin{equation}
\tilde{\eta}(s) = \frac{1}{s/\mu_m + 1/\eta} = \frac{\eta}{1+\tau_R s}.
\label{eq:eta_eff}
\end{equation}
At $s\to0$ (long times) $\tilde{\eta}(s)\to\eta$, recovering the bare
Newtonian dashpot; at $s\to\infty$ (short times) $\tilde{\eta}(s)\,s\to\mu_m$,
recovering the purely elastic response of the spring. Replacing $\eta$ by
$\tilde{\eta}(s)$ in the Stokes kernel (\ref{eq:stokes}) and returning to
the time domain gives the visco-elastic Stokes kernel
(Appendix~\ref{app:correspondence}):
\begin{equation}
\sigma_{xz}^{\text{BOT}} + \tau_R\frac{\partial\sigma_{xz}^{\text{BOT}}}{\partial t}
= \int_{-\infty}^{+\infty}\frac{2\eta}{\pi(x-x')^2}
  \frac{\partial u(x',t)}{\partial t}\,dx'.
\label{eq:ve_kernel}
\end{equation}
In the limit $\tau_R\to 0$, (\ref{eq:ve_kernel}) reduces to the EV
Stokes kernel (\ref{eq:stokes}).

\subsubsection{Governing equation in Fourier space}

The equilibrium equation (\ref{eq:equil}) expresses the basal traction in terms
of the elastic crust displacement:
\begin{equation}
\sigma_{xz}^{\text{BOT}} = -\frac{EH}{1-\nu^2}\frac{\partial^2 u}{\partial x^2}.
\label{eq:sigma_from_equil}
\end{equation}
Substituting (\ref{eq:sigma_from_equil}) and its time derivative into
(\ref{eq:ve_kernel}), then applying the spatial Fourier transform
using (\ref{eq:stokes_FT}):
\begin{equation}
\frac{EH}{1-\nu^2}k^2\hat{u}
+ \tau_R\frac{EH}{1-\nu^2}k^2\frac{\partial\hat{u}}{\partial t}
= -2\eta|k|\frac{\partial\hat{u}}{\partial t}.
\label{eq:maxwell_fourier_pre}
\end{equation}
Dividing by $2\eta$, using (\ref{eq:a}), and collecting
$\partial\hat{u}/\partial t$ terms gives the visco-elastic analogue of
(\ref{eq:fourier}):
\begin{equation}
\frac{\partial\hat{u}}{\partial t}
= -\left(\frac{a|k|}{1+\tau_R a|k|}\right)\hat{u}.
\label{eq:maxwell_fourier}
\end{equation}
The denominator $1+\tau_R a|k|$ acts as a low-pass filter on the relaxation
rate, slowing the decay of high-wavenumber (near-field) features relative to
the EV case.

\subsection{Step response and asymptotic behavior}
\label{sec:eve_step}

Equation (\ref{eq:maxwell_fourier}) is separable in $k$, yielding a decaying
exponential for each wavenumber:
\begin{equation}
\hat{u}(k,t) = \hat{u}(k,0)\exp\!\left(-\frac{a|k|}{1+\tau_R a|k|}\,t\right).
\label{eq:maxwell_fourier_sol}
\end{equation}
Applying the Heaviside initial condition (\ref{eq:IC_FT}):
\begin{equation}
\hat{u}(k,t) = b\pi\delta(k)
  + \frac{ib}{k}\exp\!\left(-\frac{a|k|}{1+\tau_R a|k|}\,t\right).
\label{eq:maxwell_IC_evolved}
\end{equation}
Taking the inverse Fourier transform and following the same parity argument as
in Section~\ref{sec:step}, the EVE step displacement requires numerical
integration:
\begin{equation}
u^{\text{EVE}}(x,t) = \frac{b}{2}
  - \frac{b}{\pi}\int_0^{+\infty}
    \frac{\sin(kx)}{k}\exp\!\left(-\frac{akt}{1+\tau_R ak}\right)dk,
\label{eq:maxwell_step}
\end{equation}
since the rational exponent prevents the closed-form Dirichlet evaluation of
Appendix~\ref{app:step}. The surface velocity follows by differentiating
(\ref{eq:maxwell_step}) with respect to $t$:
\begin{equation}
v^{\text{EVE}}(x,t) = \frac{b a}{\pi}\int_0^{+\infty}
  \frac{\sin(kx)}{1+\tau_R ak}
  \exp\!\left(-\frac{akt}{1+\tau_R ak}\right)dk.
\label{eq:maxwell_vel}
\end{equation}
Both integrals are evaluated numerically (Appendix~\ref{app:numerics}).

The exponent in (\ref{eq:maxwell_step}) reveals two physically distinct regimes.
In the long-wavelength limit ($k\to 0$ or $t\to\infty$), $1+\tau_R ak \approx 1$
and the exponential reduces to $\exp(-akt)$, recovering the EV
Dirichlet integral of Section~\ref{sec:step}. The elastic memory of the
asthenosphere fades at large distances and late times, and the model converges
to the Newtonian viscous limit. In the short-wavelength limit ($k\to\infty$ or
$t\to 0$), $1+\tau_R ak \approx \tau_R ak$ and the exponent saturates:
\begin{equation}
\lim_{k\to\infty}\left(-\frac{akt}{1+\tau_R ak}\right) = -\frac{t}{\tau_R}.
\label{eq:maxwell_saturation}
\end{equation}
At early times $t\ll\tau_R$, the exponential approaches unity and no near-field
deformation occurs outside the rift interface. For $t\sim\tau_R$, near-field
deformation decays exponentially at the fixed rate $1/\tau_R$, representing
pure Maxwell relaxation of the mantle independent of the elastic crust.
Together, (\ref{eq:maxwell_vel}) and (\ref{eq:maxwell_saturation}) demonstrate
that the EVE model introduces a short-time, near-field stiffening relative
to the EV model, and that the EVE step response converges to the EV one
at $t\gg\tau_R$.

The saturated exponent also gives the near-axis velocity in closed form.
For $k \gg 1/(a\tau_R)$, the integrand of (\ref{eq:maxwell_vel}) reduces
to $\sin(kx)\,e^{-t/\tau_R}/(\tau_R k)$, and these wavenumbers control the
velocity as $x \to 0$. The Dirichlet integral of Section~\ref{sec:step}
then gives
\begin{equation}
v^{\text{EVE}}(x\to0^+,\,t) = \frac{b}{2\tau_R}\,e^{-t/\tau_R},
\label{eq:maxwell_plateau}
\end{equation}
a finite plateau that decays at the Maxwell rate, whereas the EV velocity
vanishes at the axis. The near-field signal of the EVE model is therefore
an elastic rebound of the half-space, paced by $\tau_R$ rather than by the
diffusion front $a t$.

Figure~\ref{fig:maxwell_halfspace_step} shows the EVE STEP response in
the same layout as its EV counterpart (Figure~\ref{fig:viscous_halfspace_step}):
same velocity unit $b/\tau_R$, colour scale, windows, vertical guide
$x = a\tau$ and phases, except that its first row shows the limit
$\tau\to0^+$ rather than $\tau = 0.001\,\tau_R$. Any difference between
the two figures therefore reflects the rheology alone. Unlike the EV
case, the surface velocity is finite immediately after the intrusion and
reaches $b/(2\tau_R)$ at the axis, the near-field plateau of
Equation~(\ref{eq:maxwell_plateau}). Peak
surface velocity is largest at $\tau\to0^+$ and falls with increasing
$\tau$, and the underlying flow lobe shrinks and thins in step, the same
qualitative decay as the EV STEP field but paced by $\tau_R$ rather than
by the diffusion front.

\begin{figure*}
  \centering
  \includegraphics[width=\textwidth]{fig/fig_halfspace_step_maxwell_contour.png}
  \caption{EVE STEP response to a single dike intrusion of opening $b$,
    with the same layout, units, colour scale and parameters as the EV
    STEP response (Figure~\ref{fig:viscous_halfspace_step}) and
    $\tau_R=\eta/\mu_m$ from the shared-modulus closure
    (Equation~\ref{eq:tau_R_shared}). Rows show
    $\tau/\tau_R \in \{0^+,\,0.25,\,0.5,\,1,\,2,\,3\}$. The first row is
    the limit $\tau\to0^+$, evaluated in closed form: unlike the EV case,
    the surface velocity there is finite, reaching $b/(2\tau_R)$ at the
    axis (Equation~\ref{eq:maxwell_plateau}).
    Left column: surface velocity $v^{\text{EVE}}$
    (Equation~\ref{eq:maxwell_vel}; solid, left axis, in units of
    $b/\tau_R$) and displacement $u^{\text{EVE}}$
    (Equation~\ref{eq:maxwell_step}; dashed, right axis, in units of $b$),
    counted from the instant after the intrusion. The dashed vertical
    line marks $x=a\tau$, the EV velocity-peak position, for direct
    comparison with Figure~\ref{fig:viscous_halfspace_step}.
    Middle column: speed $|v|$ in the half-space (filled contours,
    logarithmic scale, in units of $b/\tau_R$), computed with
    Equation~(\ref{eq:halfspace_kernel}) over $|x|\le 30H$ and
    $z\le 30H$. Arrows show the flow direction only and have a fixed
    length.
    Right column: the same field enlarged to $|x|\le H$, $z\le H$ beneath
    the rift axis, with the same colour scale. Speeds outside the colour
    bar range $[0.01,\,1]$ are clipped to its end colours.}
  \label{fig:maxwell_halfspace_step}
\end{figure*}

\subsection{EVE Recurrent Steady State}
\label{sec:maxwell_css}

\subsubsection{Infinite sum and Fourier periodicity condition}

The EVE RSS is the superposition of step-response velocities from all
preceding intrusions, by the same argument as in Section~\ref{sec:css}:
\begin{equation}
  v_{\text{RSS}}^{\text{EVE}}(x,\tau)
  = \sum_{j=0}^{\infty} v^{\text{STEP,EVE}}(x,\,\tau + jT).
\end{equation}
Writing $v^{\text{STEP,EVE}}$ in Fourier space
(Section~\ref{sec:maxwell_model}, Equation~\ref{eq:maxwell_vel}) as
$\frac{b a}{\pi}\int_0^\infty \frac{\sin(kx)}{1+\tau_R ak}\,e^{-\alpha(k)(\tau+jT)}\,dk$,
with the effective Maxwell decay rate
\begin{equation}
  \alpha(k) := \frac{ak}{1+\tau_R ak},
  \label{eq:alpha_k}
\end{equation}
each past intrusion contributes a factor $e^{-\alpha(k)jT}$.  The series
converges for all $k>0$ and the geometric sum
$\sum_{j=0}^{\infty}e^{-\alpha(k)jT} = (1-e^{-\alpha(k)T})^{-1}$ can be
performed inside the Fourier integral.  The same result follows from requiring each Fourier mode to repeat from
one cycle to the next, apart from the offset added by each intrusion,
$\hat{u}(k,0^+) = \hat{u}(k,T^-) + ib/k$.  Together with the decay
(\ref{eq:maxwell_fourier}) within the cycle, this condition has the unique
solution
\begin{equation}
  \hat{u}(k,\tau) = \frac{ib/k}{1-e^{-\alpha(k)T}}\,e^{-\alpha(k)\tau},
  \qquad \tau\in[0,T).
  \label{eq:maxwell_css_fourier}
\end{equation}

\subsubsection{Physical-space representation, conservation, and limits}

Inverting (\ref{eq:maxwell_css_fourier}) gives the EVE RSS velocity
and displacement:
\begin{align}
  v_{\text{RSS}}^{\text{EVE}}(x,\tau)
  &= \frac{b a}{\pi}\int_0^{\infty}
     \frac{\sin(kx)}{1+\tau_R ak}\,
     \frac{e^{-\alpha(k)\tau}}{1-e^{-\alpha(k)T}}\,dk,
  \label{eq:maxwell_css_v} \\
  u_{\text{RSS}}^{\text{EVE}}(x,\tau)
  &= \frac{b}{\pi}\int_0^{\infty}
     \frac{\sin(kx)}{k}\,
     \frac{1-e^{-\alpha(k)\tau}}{1-e^{-\alpha(k)T}}\,dk.
  \label{eq:maxwell_css_u}
\end{align}
Integrating the velocity over one full period, using
$\int_0^T e^{-\alpha\tau}/(1-e^{-\alpha T})\,d\tau = 1/\alpha$ and the
algebraic identity $a/[(1+\tau_R ak)\,\alpha(k)] = 1/k$, gives
\begin{equation}
  \int_0^T v_{\text{RSS}}^{\text{EVE}}(x,\tau)\,d\tau
  = \frac{b}{\pi}\int_0^{\infty}
    \frac{\sin(kx)}{k}\,dk = \frac{b}{2},
\end{equation}
where the last step is the Dirichlet integral. Half the dike opening,
$b/2$, is recovered every cycle, independent of $\tau_R$, matching the EV
RSS result (\ref{eq:conservation}) and the late-time limit $b/2$ of the
single-event EVE displacement (\ref{eq:maxwell_step}). In the limit
$\tau_R\to 0$, $\alpha(k)\to ak$ and $(1+\tau_R ak)^{-1}\to 1$, and
(\ref{eq:maxwell_css_v}) reduces to the EV RSS Fourier integral whose
closed form is the Digamma series of Section~\ref{sec:css}.  For
$\tau_R > 0$, the saturating denominator $1+\tau_R ak$ prevents this
reduction and (\ref{eq:maxwell_css_v}) is evaluated numerically
(Appendix~\ref{app:numerics}).

The interior flow of the EVE RSS again follows from
Equation~(\ref{eq:halfspace_kernel}), with $V$ set to
$v_{\text{RSS}}^{\text{EVE}}$ (\ref{eq:maxwell_css_v}).
Figure~\ref{fig:maxwell_halfspace_css} shows it in the same layout,
normalisation and phases as the EV RSS
(Figure~\ref{fig:viscous_halfspace_css}), for the same $T = 5\tau_R$,
except that its first row shows the limit $\tau\to0^+$. Far from the axis
($|x| \gg aT$), each flank moves at $V/2$ throughout the cycle, as in the
EV RSS. Near the axis, the cycle starts from a finite velocity jump at the
source. Each past intrusion contributes the plateau
(\ref{eq:maxwell_plateau}) at its own age $\tau + jT$, and summing these
contributions gives
\begin{equation}
v_{\text{RSS}}^{\text{EVE}}(x\to0^+,\,\tau)
= \frac{b}{2\tau_R}\,\frac{e^{-\tau/\tau_R}}{1-e^{-T/\tau_R}},
\label{eq:maxwell_rss_jump}
\end{equation}
so the velocity jumps by twice this value across the axis for any
$\tau_R>0$, whereas the EV RSS vanishes continuously there. Immediately
after an intrusion, the near-axis velocity is about
$2.5\,V$ (Equation~\ref{eq:maxwell_rss_jump}), the elastic rebound of the
most recent intrusion together with the remnants of earlier ones. This
jump decays at the rate $1/\tau_R$, to about $0.9\,V$ at
$\tau = 0.2\,T = \tau_R$ and $0.3\,V$ at $\tau = 0.4\,T$. By the end of
the cycle it has all but vanished, and the profile has smoothed into a
monotonic ramp from $-V/2$ to $+V/2$ over a distance of order $aT$, as in
the EV RSS. The interior field follows the same pattern as in the EV
case. Upwelling beneath the axis and outward flow beneath the plate persist
at every phase, and the speed over most of the half-space stays between
roughly $0.1\,V$ and $V/2$ (middle column). A low-velocity region opens
beneath the axis after each intrusion and grows until the next one (right
column).

\begin{figure*}
  \centering
  \includegraphics[width=\textwidth]{fig/fig_halfspace_css_maxwell_contour.png}
  \caption{EVE RSS for $T=5\tau_R$, in the same layout, units and colour
    scale as the EV RSS (Figure~\ref{fig:viscous_halfspace_css}), with
    $\tau_R=\eta/\mu_m$ from the shared-modulus closure
    (Equation~\ref{eq:tau_R_shared}), at six phases
    $\tau/T \in \{0^+,\,0.2,\,0.4,\,0.6,\,0.8,\,0.99\}$, one per row.
    The first row is the limit $\tau\to0^+$, computed from its exact
    form: the displacement since the most recent intrusion is zero, and
    the surface velocity at the axis,
    $(b/2\tau_R)/(1-e^{-T/\tau_R})\approx2.5\,V$
    (Equation~\ref{eq:maxwell_rss_jump}), lies beyond the plotted range
    and is clipped at $2\,V$.
    Left column: surface velocity $v_{\text{RSS}}^{\text{EVE}}$
    (Equation~\ref{eq:maxwell_css_v}; solid, left axis, in units of
    $V=b/T$) and displacement accumulated since the most recent intrusion
    $u_{\text{RSS}}^{\text{EVE}}$ (Equation~\ref{eq:maxwell_css_u}; dashed,
    right axis, in units of $b$).  The dashed vertical line marks
    $x=a\tau$, the EV velocity-peak position of the most recent intrusion
    alone, for direct comparison with
    Figure~\ref{fig:viscous_halfspace_css}.
    Middle column: speed $|v|$ in the half-space (filled contours,
    logarithmic scale, in units of $V$), computed with
    Equation~(\ref{eq:halfspace_kernel}) over $|x|\le 4.5\,aT$ and
    $z\le 4.5\,aT$.  Arrows show the flow direction only and have a fixed
    length.
    Right column: the same field enlarged to $|x|\le 0.15\,aT$,
    $z\le 0.15\,aT$ beneath the rift axis, with the same colour scale.
    Speeds outside the colour bar range $[0.01,\,1]$ are clipped to its
    end colours.  Parameters as in Figure~\ref{fig:viscous_halfspace_step},
    giving $aT \approx 67$~km.}
  \label{fig:maxwell_halfspace_css}
\end{figure*}

\subsection{Constraining $\tau_R$ without an extra free parameter}
\label{sec:tauR_shared}

Section~\ref{sec:results} fits the EVE RSS to real data alongside its
EV counterpart.  A free $\tau_R$ alongside $(a,\,b)$ would let the extra
parameter absorb noise rather than test whether the EVE model improves the
fit: a model with one more free parameter always reaches at least as low a
$\chi^2$, whether or not that parameter is physically warranted.  We
instead derive $\tau_R$ from $a$, so that the EVE RSS carries exactly as
many free parameters as the EV RSS it is compared against.

Doing so requires the shear modulus $\mu_m$ of the half-space, left
general in Equation~(\ref{eq:eta_eff}) so far. By definition,
$\tau_R := \eta/\mu_m$ is a material property of the half-space alone.
The EV half-space is purely viscous and has no elastic modulus of its own:
its decay rate $a|k|$ (\ref{eq:fourier}) depends on wavenumber and is not
a material relaxation rate. The analogous timescale of the coupled
plate--half-space \emph{system} is the same viscosity-over-modulus ratio,
with the shear modulus of the plate in place of that of the half-space:
\begin{equation}
\tau_{\text{EV}} := \frac{\eta}{\mu_{\text{plate}}}, \qquad
\mu_{\text{plate}} := \frac{E}{2(1+\nu)}.
\label{eq:tau_EV}
\end{equation}
$\tau_R$ and $\tau_{\text{EV}}$ therefore differ only in which shear
modulus, $\mu_m$ or $\mu_{\text{plate}}$, is paired with the half-space
viscosity $\eta$.

We adopt the simplest closure available without additional data: the
half-space shares the same shear modulus as the elastic plate it
underlies,
\begin{equation}
\mu_m = \mu_{\text{plate}} := \frac{E}{2(1+\nu)}.
\label{eq:mu_shared}
\end{equation}
This closure makes the two relaxation times identical by construction,
with no free parameter left to tune the agreement:
\begin{equation}
\tau_R = \frac{\eta}{\mu_m} = \frac{\eta}{\mu_{\text{plate}}} =
  \tau_{\text{EV}}.
\end{equation}
Combined with $\eta = EH/[2a(1-\nu^2)]$, the inversion of the characteristic
velocity (\ref{eq:a}), this gives
\begin{equation}
\tau_R = \tau_{\text{EV}} = \frac{H}{a(1-\nu)},
\label{eq:tau_R_shared}
\end{equation}
a relaxation time set entirely by the fitted characteristic velocity $a$
and the canonical parameters $H$ and $\nu$ (Section~\ref{sec:disc_rheology}),
with no independent estimate of $\mu_m$ or $\eta$.  The shared-modulus
closure is a simplifying assumption, not a physical necessity: the mantle
shear modulus need not equal that of the crust.  It is, however, the
natural default in the absence of independent constraints, and it keeps
the EV-versus-EVE comparison of Section~\ref{sec:results} an
equal-parameter test of functional form.

% ====================================================================
\section{Applications to magmatic rift systems}
\label{sec:results}

We test the EV and EVE Recurrent Steady States
(Sections~\ref{sec:css} and~\ref{sec:maxwell_css}) against post-diking
geodetic records from two magmatic rifts. The first is the 1975--1984
Krafla rifting episode in NE Iceland, sampled by two rift-perpendicular
displacement profiles \citep{foulger1992post,hofton1996post}. The second
is the 1978 Asal--Ghoubbet episode in the Afar Triangle, sampled by a
25-year baseline time series and a rift-perpendicular velocity profile
\citep{vigny2007twenty}. The two sites differ in rift opening rate, in the
length of their records and in the type of data, so together they test
whether a single functional form describes inter-rifting deformation in
contrasting settings.

Both records cover only part of one cycle. They span 12--17~yr after the
onset of the Krafla episode and 1--25~yr after the Asal event, and in
neither case has a later intrusion been documented. We set the recurrence
interval through the far-field conservation of the RSS
(Equation~\ref{eq:conservation}). Each flank advances by $b/2$ per cycle,
so the two flanks separate by the full dike opening $b$ every $T$, and the
rift opening rate, the rate at which the two flanks separate, is
$V = b/T$, with each flank moving at $\pm V/2$ far from the axis. We take
$V$ from an independent, long-term estimate of the rift opening rate, from
global plate motion models or regional geodetic studies, and derive the
recurrence interval from the fitted opening as $T = b/V$.
This leaves two free parameters, the characteristic velocity $a$ and the
dike opening $b$, plus a reference-frame offset where the data require
one. The EVE RSS carries $\tau_R$ derived from $a$ through the
shared-modulus closure (Equation~\ref{eq:tau_R_shared}), so each EVE fit
has exactly as many free parameters as its EV counterpart. Although the
records are short, the RSS still contains the contribution of all earlier
intrusions, including the far-field plate motion they sustain.

At each site we minimise the misfit
\begin{equation}
\chi^2 = \sum_{i=1}^{N}
  \left(\frac{u_i^{\mathrm{obs}} - u_i^{\mathrm{pred}}}{\sigma_i}\right)^2,
\label{eq:chi2}
\end{equation}
where $u_i^{\mathrm{obs}}$ and $u_i^{\mathrm{pred}}$ are the observed and
predicted displacement, baseline change or velocity at data point $i$,
$\sigma_i$ is its $1\sigma$ uncertainty, and $N$ is the number of
observations. We report the reduced statistic $\chi^2_r = \chi^2/(N-p)$,
where $p$ is the number of free parameters, so that $\chi^2_r \approx 1$
indicates a fit within measurement noise. We compare the EV and EVE fits
with the Akaike information criterion (AIC). Because both models carry
the same number of free parameters, the AIC difference reduces to the
difference in $\chi^2$, and any preference reflects the shape of the
predicted deformation alone.

\begin{table*}
\centering
\caption{Best-fit RSS parameters for the joint Krafla and Asal fits, with
  $1\sigma$ ranges in brackets ($\Delta\chi^2 \le 2.30$ projections, two
  parameters). The recurrence interval is derived as $T=b/V$
  ($V=1.8$~cm~yr$^{-1}$ at Krafla, 13~mm~yr$^{-1}$ at Asal) and $\tau_R$
  from $a$ through (\ref{eq:tau_R_shared}); $c$ is the reference-frame
  offset of the \citet{hofton1996post} profile, and $k$ the number of free
  parameters. At Asal, $\chi^2/N$ is 1.37 (EV) and 1.62 (EVE) on the 33
  time-series points and 21.3 and 25.4 on the 15 profile points.
  $\Delta$AIC is relative to the EV RSS at the same site.}
\label{tab:fits}
\begin{tabular}{llcccccccc}
\hline
Site & Model & $a$ ($10^{-5}$~m~s$^{-1}$) & $b$ (m) & $T$ (yr) & $\tau_R$ (yr) & $c$ (cm) & $k$ & $\chi^2_r$ & $\Delta$AIC \\
\hline
Krafla ($N=37$) & EV RSS  & 4.99 [4.21, 5.82] & 1.71 [1.60, 1.80] & 95 [89, 100] & --              & 6.5 & 3 & 1.61 & --      \\
                & EVE RSS & 11.9 [10.8, 13.1] & 1.60 [1.51, 1.69] & 89 [84, 94]  & 3.6 [3.2, 3.9]  & 6.6 & 3 & 1.82 & $+7.1$  \\
\hline
Asal ($N=48$)   & EV RSS  & 1.00 [0.97, 1.04] & 0.673 [0.666, 0.681] & 52 [51, 52] & --             & --  & 2 & 7.93 & --      \\
                & EVE RSS & 6.95 [6.77, 7.14] & 0.642 [0.637, 0.648] & 49 [49, 50] & 6.1 [5.9, 6.2] & --  & 2 & 9.46 & $+70.3$ \\
\hline
\end{tabular}
\end{table*}

\subsection{Krafla rift, Iceland}
\label{sec:krafla}

The Krafla rifting episode began in December 1975 and lasted about a
decade, with some 20 dike intrusions along the Krafla fissure swarm
\citep{hofton1996post}. Its aftermath was recorded by GPS campaigns in
1987, 1990 and 1992. We digitized two rift-perpendicular displacement
profiles from these campaigns. The first comprises 20 stations of
\citet{foulger1992post} (their Figure~2b), differenced between 1987 and
1990. The second comprises 17 stations of \citet{hofton1996post} (their
Figure~5, the same profile they compare with the diffusion model in their
Figure~16), differenced between 1987 and 1992. Each point carries its
published $1\sigma$ error bar. The two profiles also differ in their
spatial reference. \citet{foulger1992post} tie their displacements to
near-rift stations that showed no significant deformation, whereas
\citet{hofton1996post} state that the zero level of their profile is
arbitrary. We therefore add a constant offset $c$ to the model for the
\citet{hofton1996post} profile alone and solve for it during the fit.

A joint fit first requires a common time origin for the two profiles.
Our RSS represents each episode as a single intrusion of opening $b$, and
it predicts each profile from the elapsed times $(\tau_i,\tau_f)$ since
that intrusion, so both datasets must measure these times from the same
instant. Neither published study provides one directly: both measure time
from an origin that differs from the onset of the episode.
\citet{foulger1992post} treat their 1987--1990 displacements as an
instantaneous velocity 11~yr after the dike intrusions, and
\citet{hofton1996post} treat their 1987--1992 displacements as one 13~yr
after, so both place the intrusion in 1979. \citet{hofton1996post}
likewise inject all dikes of their visco-elastic model simultaneously in
1979, and report that spreading the injections over the real nine-year
sequence changes the 1987--1992 field negligibly. We adopt the onset of
the episode, $t_0 = 1975$, for both datasets. The half-space begins to relax as
soon as the first dike opens the elastic plate, in December 1975. An
origin in 1979 would treat the stresses imposed by the intrusions of the
preceding four years as if their relaxation had not yet started. Because
the opening was in fact spread over the episode, $t_0 = 1975$ gives the
longest elapsed time compatible with the record, and the times since
intrusion used below are upper bounds. The two epoch pairs then sit at
$\tau = 12$--$15$~yr and $12$--$17$~yr after the intrusion.

With this common origin, we fit both profiles jointly, because they
record the same episode.  A single $(a,b)$ is shared by the two profiles,
each is evaluated at its own $(\tau_i,\tau_f)$, and the pooled $\chi^2$
(\ref{eq:chi2}) is minimized over the $N=37$ stations combined, with $c$
as a third free parameter. The recurrence interval follows from $T = b/V$ with
$V=1.8$~cm~yr$^{-1}$, the NUVEL-1A full spreading rate
for north Iceland \citep{demets1994effect}, which \citet{hofton1996post}
adopt as the time-averaged opening rate across the region.
\citet{foulger1992post} used the 1.9~cm~yr$^{-1}$ predicted by the earlier
NUVEL-1 model, which the NUVEL-1A revision of the geomagnetic time scale
superseded.

Figure~\ref{fig:krafla_v_constrained} and Table~\ref{tab:fits} summarize
the result.  Both models reproduce the near-rift step and the far-field
plateau of the two profiles (Figure~\ref{fig:krafla_v_constrained}a,b)
with $\chi^2_r < 2$, and the data favour the EV RSS at the same parameter
count ($\Delta\mathrm{AIC}=+7.1$).

Both $\chi^2$ landscapes are closed and single-minimum within the
$3\sigma$ contour (Figure~\ref{fig:krafla_v_constrained}c,d), so $a$ and
$b$ are resolved independently despite the short observation window
(Table~\ref{tab:fits}).  In the EVE landscape $a$ and $b$ are positively
correlated, while in the EV landscape they are nearly independent.  Section~\ref{sec:disc_eve}
discusses why the two rheologies are hard to separate at Krafla, and
Section~\ref{sec:disc_rheology} the implied viscosity and recurrence
interval.

\begin{figure*}
  \centering
  \includegraphics[width=\textwidth]{fig/fig_krafla_v_constrained_combined.png}
  \caption{Joint EV and EVE RSS fits to both Krafla datasets, with $T$
    derived from the fitted dike opening $b$ through the opening-rate
    relation $T=b/V$ at $V=1.8$~cm~yr$^{-1}$ and time measured from
    $t_0=1975$. (a) \citet{foulger1992post} ($\tau_i=12$~yr,
    $\tau_f=15$~yr). (b) \citet{hofton1996post} ($\tau_i=12$~yr,
    $\tau_f=17$~yr), plotted at its own uncorrected displacements, with the
    reference-frame offset absorbed into each curve rather than removed from
    the data. Curves: EV RSS (purple) and EVE RSS (navy), each with one
    $(a,b)$ shared by both panels; error bars are $1\sigma$. (c, d) Joint
    $\chi^2$ landscapes in $(a,b)$ for the EV and EVE RSS, shown as
    $\Delta\chi^2=\chi^2-\chi^2_{\min}$, with the reference-frame offset
    minimised at every grid point. Contours mark the 68.3, 95.4 and 99.7 per
    cent confidence levels for two parameters ($\Delta\chi^2 = 2.30$, $6.17$,
    $11.83$); the cross marks the best fit. The right axis gives the derived
    $T=b/V$. See Table~\ref{tab:fits} for goodness-of-fit
    statistics.}
  \label{fig:krafla_v_constrained}
\end{figure*}

\subsection{Asal rift, Djibouti}
\label{sec:asal}

The November 1978 Asal--Ghoubbet rifting episode in the Afar Triangle opened
approximately 2~m of rift-perpendicular extension over several days \citep{ledain1979seismotectonic,ruegg1984seismotectonic}.  \citet{vigny2007twenty}
tracked the EP00--DF00 baseline, a rift-perpendicular geodetic tie crossing
the Asal rift, from 1979 to 2003.  The baseline lengthened by
$\approx0.5$~m, rapidly at first and much more slowly after the
mid-1980s.  \citet{vigny2007twenty} also measured rift-perpendicular GPS
velocities along a transect crossing the rift, at a mean epoch of
$\approx19$~yr after the event.  We digitized all 33 time-series points
from their Figure~3 and all 15 profile points from their velocity
transect, and fit both jointly with the EV and EVE RSS forward models of
Sections~\ref{sec:css} and~\ref{sec:maxwell_css}.  The predicted baseline
change is $\Delta L = 2u(x_0,\tau)$ at $x_0 = 2000$~m, half the EP00--DF00
rift-perpendicular separation, with the rifting onset fixed at
$t_0 = 1978.85$~yr \citep{vigny2007twenty}.  The predicted velocity profile
is evaluated at the transect epoch, with the model centre placed at
$x=+1$~km northeast of the topographic rift axis, following
\citet{vigny2007twenty}.  Because the \citeauthor{vigny2007twenty} time
series is referenced to the first measurement epoch
($\tau_{\mathrm{ref}} \approx 1.1$~yr after the rifting event) rather than to
the rifting onset, we subtract from each model curve its own prediction at
$\tau_{\mathrm{ref}}$, anchoring $\Delta L = 0$ at the reference epoch without
introducing a free parameter.

As at Krafla, the recurrence interval follows from the opening-rate
relation $T = b/V$.  We take
$V = 13$~mm~yr$^{-1}$, the long-term Arabia--Somalia rate reported by
\citet{vigny2007twenty}.  We do not use the $\approx16$~mm~yr$^{-1}$ opening
they measure across the Asal rift itself: they attribute its excess over
the long-term rate to the post-1978 transient, which the RSS already produces,
so adopting it as $V$ would count the transient twice.  The velocity
profile is expressed in the Somalia reference frame, in which the SW flank
is nearly fixed and the NE flank moves at the full rate.  Because the
model's far-field asymptotes are $\pm V/2$, the frame shift is $V/2$ by
construction rather than a fitted offset.  No second intrusion has been
documented at Asal--Ghoubbet since 1978: \citet{smittarello2016transient}
analyse over 40~yr of geodetic monitoring (1972--2016) as the relaxation of
that single event.  We therefore require $T \geq T_{\min} = 37.2$~yr, the
1978--2016 span, which under $T = b/V$ becomes the lower bound
$b \geq V\,T_{\min} \approx 0.48$~m.  Only $(a,b)$ remain free, and the EVE
variant carries $\tau_R$ derived from $a$ through (\ref{eq:tau_R_shared}).

Figure~\ref{fig:asal_v_constrained} and Table~\ref{tab:fits} summarize
the result.  Neither fit sits on the $T_{\min}$ floor.  Both
models reproduce the time series closely, with $\chi^2/N = 1.37$ for the
EV RSS and $1.62$ for the EVE RSS over its 33 points.  The elevated joint
$\chi^2_r$ comes from the profile ($\chi^2/N = 21.3$ and $25.4$ over its 15
points), where the observed SW--NE asymmetry exceeds what a single
symmetric transient can produce.  The EVE profile also shows a velocity
jump of $\approx4.6$~mm~yr$^{-1}$ at the model centre
(Figure~\ref{fig:asal_v_constrained}b), the RSS jump of
Equation~(\ref{eq:maxwell_rss_jump}) at $\tau \approx 19$~yr.  The EV RSS
fits better at equal parameter count, with $\Delta\mathrm{AIC} = +70.3$ in
its favour.  Section~\ref{sec:disc_eve} discusses why the Asal record
separates the two rheologies, and Section~\ref{sec:disc_rheology} the
implied viscosity and the field observations that the fits do not
reproduce.

Figure~\ref{fig:asal_v_constrained}c,d shows the joint (time series +
profile) $\chi^2$ landscapes for both rheologies.  The combined dataset
confines $(a,b)$ to a narrow, anti-correlated ellipse in both cases, well
above the $T_{\min}$ floor, with $1\sigma$ ranges of a few per cent
(Table~\ref{tab:fits}).

\begin{figure*}
  \centering
  \includegraphics[width=\textwidth]{fig/fig_asal_v_constrained_combined.png}
  \caption{Joint EV and EVE RSS fits to the Asal data of
    \citet{vigny2007twenty}, with $T$ derived from the fitted dike opening
    $b$ through the opening-rate relation $T=b/V$ at $V=13$~mm~yr$^{-1}$.
    (a) EP00--DF00 baseline change (mm), filled circles with $1\sigma$
    error bars; model curves are referenced to the first measurement epoch
    ($\tau_{\mathrm{ref}} \approx 1.1$~yr). (b) Rift-perpendicular
    velocity profile at $\tau\approx19$~yr in the Somalia reference frame
    (mm~yr$^{-1}$); model curves include the frame shift $V/2$. The dotted
    line marks the topographic rift axis and the dashed line the model
    centre at $x=+1$~km. Curves: EV RSS (purple) and EVE RSS (navy), each
    with one $(a,b)$ shared by both panels. (c, d) Joint $\chi^2$
    landscapes in $(a,b)$ for the EV and EVE RSS, shown as
    $\Delta\chi^2=\chi^2-\chi^2_{\min}$. Contours mark the 68.3, 95.4 and
    99.7 per cent confidence levels for two parameters
    ($\Delta\chi^2 = 2.30$, $6.17$, $11.83$); the cross marks the best fit.
    The right axis gives the derived $T=b/V$. See Table~\ref{tab:fits}
    for goodness-of-fit statistics.}
  \label{fig:asal_v_constrained}
\end{figure*}

% ====================================================================
\section{Discussion}
\label{sec:discussion}

% ── §4.1 ─────────────────────────────────────────────────────────────────────
\subsection{Inter-rifting deformation as the memory of discrete intrusions}
\label{sec:disc_bcs}

Kinematic models of inter-rifting deformation need a free parameter that
the mechanics of a magmatic rift does not supply.  Backslip models
\citep{savage1978asthenosphere} and their dislocation equivalents
\citep{lafemina2005geodetic,konca2018slip} reproduce observed velocity
gradients through a locking depth or a creep-zone geometry, which carries
no independent constraint where the rheological transition is gradational.
Numerical visco-elastic simulations face the same obstacle at their lateral
boundaries: these must either force the far-field velocity to zero, in
contradiction with global plate motion models \citep{altamimi2016itrf2014},
or prescribe the divergence rate that the model is meant to predict
\citep{hofton1995horizontal,hofton1996post}.

The RSS removes both requirements, because it treats inter-rifting
deformation as the superposed relaxation of all past intrusions rather
than as a separate process.  Far-field plate divergence follows from the
conservation of the opening over each cycle
(Equation~\ref{eq:conservation}), not from a boundary condition, and the
only length scale is the deformation half-width
$aT = EHT/[2\eta(1-\nu^2)]$, set by the balance between elastic stresses in
the plate and viscous drag at its base.  A station that records a steady
velocity far from the rift therefore does not reveal continuous creep at
depth: it records the cumulative relaxation of intrusions that occurred
many cycles earlier.

This memory sets the conditions under which the RSS applies.  Truncating
the history after $N$ intrusions leaves a velocity deficit of about
$(V/\pi)\arctan[|x|/(NaT)]$ (Section~\ref{sec:css}), so the RSS is
reproduced to within 5 per cent of the far-field velocity $V/2$ only where
$N \gtrsim 13\,|x|/(aT)$.  The number of intrusions a rift must have hosted
thus grows linearly with distance from its axis, a direct consequence of
the slowly decaying $1/x$ tail of the step response
(Section~\ref{sec:step_selfsimilar}).  The farthest stations used here lie
at $0.9\,aT$ at Krafla and $0.6\,aT$ at Asal, so the RSS describes them
provided each rift has hosted more than about 12 and 8 intrusions, about
1100 and 400~yr of rifting at the derived recurrence intervals.  On a newly
initiated rift segment, or at stations far beyond $aT$, the explicit finite
sum (\ref{eq:finite_sum}) should replace the RSS.

% ── §4.2 ─────────────────────────────────────────────────────────────────────
\subsection{Non-local versus Fickian relaxation}
\label{sec:disc_nonlocal}

The EV model and the diffusion model of \citet{foulger1992post} differ in
the order of their governing equation, not only in their parameters.  The
Fickian model, $\partial u/\partial t = \kappa\,\partial^2 u/\partial x^2$,
damps each Fourier mode at the rate $\kappa k^2$; the EV model
(Equation~\ref{eq:fourier}) damps it at $a|k|$, the Fourier signature of the
half-Laplacian.  This operator arises from the $1/(x-x')^2$ Stokes kernel
(\ref{eq:stokes}): the half-space transmits stress instantaneously, so
every point of the plate feels an intrusion from the first instant, and no
diffusive front is needed to reach it.  The step velocity is therefore a
Lorentzian with a $1/x$ tail (\ref{eq:stepv_dimless}), against the Gaussian
decay of the Fickian velocity (\ref{eq:foulger_v_dimless},
Figure~\ref{fig:step_selfsimilar}b).

The two solutions are nevertheless exactly related.  By Bochner
subordination \citep{metzler2000random}, the EV decay factor is an average
of Fickian decay factors over an effective diffusion time $s$,
\begin{equation}
e^{-a|k|t} = \int_0^{\infty} e^{-k^2 s}\, g_t(s)\,ds, \qquad
g_t(s) = \frac{at}{2\sqrt{\pi}}\, s^{-3/2}\,
  \exp\!\left(-\frac{a^2t^2}{4s}\right),
\label{eq:subordination}
\end{equation}
so the EV step displacement (\ref{eq:step}) equals the Fickian profile
$(b/2)\,\mathrm{erfc}[x/(2\sqrt{s})]$ averaged over $s$ with the weight
$g_t(s)$.  We checked this identity numerically to a relative precision
better than $10^{-7}$.  The weight peaks at $s = (at)^2/6$ but decays only as
$s^{-3/2}$, so long effective diffusion times contribute substantially, and
these carry the $1/x$ tail.  The Fickian model retains a single diffusion
time, $s = \kappa t$, and discards that tail.

The two formulations leave distinct, observable signatures.  The EV
velocity peak migrates at the constant speed $a$, so a station at distance
$x$ records its peak at $t = x/a$; the Fickian peak arrives at
$t = x^2/(6\kappa)$ (Section~\ref{sec:step_selfsimilar}).  The far fields
also differ: the RSS settles onto the plate velocity $\pm V/2$, while the
Fickian profile returns to its reference level.  Krafla illustrates the
consequence.  We refit the diffusion model of
Figure~\ref{fig:hofton_diffusion} under the same conditions as the RSS:
one $(\kappa,b)$ shared by the two profiles, $t_0 = 1975$, and the offset
$c$ of the \citet{hofton1996post} profile solved during the fit.  The
refit gives $\kappa = 2.70$~m$^2$~s$^{-1}$ and $b = 2.81$~m with
$\chi^2_r = 4.90$, against 1.61 for the EV RSS at equal parameter count.  The refit diffusion
curves peak $\approx50$~km from the axis, whereas the data peak at
15--21~km, and they return to their reference level within $\sim150$~km.
Networks that span both the near-rift peak and the far-field plateau,
or time series at several distances from the axis, are therefore best
placed to discriminate between the two formulations.

% ── §4.3 ─────────────────────────────────────────────────────────────────────
\subsection{Viscous versus visco-elastic half-space: when the data can tell}
\label{sec:disc_eve}

At equal $a$, the EVE model departs from the EV model only at short
wavelengths and early times.  The elastic response of the Maxwell
half-space saturates the decay rate of modes with $k \gtrsim 1/\ell_M$,
where $\ell_M = a\tau_R$ (Equation~\ref{eq:maxwell_saturation}), and its
effect fades after a few $\tau_R$.  Its signature is a finite velocity at
the axis that decays as $e^{-t/\tau_R}$ (Equation~\ref{eq:maxwell_plateau}),
which in the RSS becomes a velocity jump across the axis
(Equation~\ref{eq:maxwell_rss_jump}).  This jump follows from treating the
dike as infinitely thin and the plate as a single layer; in reality we
expect it to be spread over a distance comparable to the plate thickness
$H$.  Under the shared-modulus closure,
$\ell_M = H/(1-\nu)$ (Equation~\ref{eq:tau_R_shared}), about 13~km for
$H = 10$~km, independently of $a$: the elastic signal is confined within a
distance of about $1.3\,H$ from the axis.

Whether a record can distinguish the two rheologies therefore depends on
two numbers: its first epoch in units of $\tau_R$, and the fraction of its
stations that lie within $\ell_M$.  At Krafla, the observations start
$3.4\,\tau_R$ after the onset, and only five of 20 and three of 17 stations
lie within $\ell_M \approx 13$~km of the axis.  Most of the elastic
contribution has decayed by then, and the data favour the EV RSS only modestly
($\Delta\mathrm{AIC} = +7.1$).  At Asal, the time series starts at
$0.2\,\tau_R$ and the whole velocity profile lies within $\ell_M$.  This
record samples the window where the two rheologies differ most, and the
comparison is decisive ($\Delta\mathrm{AIC} = +70.3$ in favour of the EV
RSS).

Two caveats qualify this preference.  First, it rests on the shared-modulus
closure, which ties $\tau_R$ and $\ell_M$ to $a$ and $H$; a mantle shear
modulus different from that of the plate would shift both scales.  Second,
the choice of rheology changes the inferred viscosity.  The EVE fits
require a larger characteristic velocity than the EV fits, by a factor of
2.4 at Krafla and 6.9 at Asal (Table~\ref{tab:fits}), which lowers $\eta$ to
$1.8\times10^{18}$ and $3.1\times10^{18}$~Pa~s at $H = 10$~km and
$E = 40$~GPa.  At Asal, this rheological uncertainty exceeds the $\chi^2$
uncertainty on $a$ by more than an order of magnitude.  Viscosities inferred
from post-diking geodesy are therefore conditional on the assumed
half-space rheology as much as on the elastic parameters.

% ── §4.4 ─────────────────────────────────────────────────────────────────────
\subsection{What the fits constrain, and where the model falls short}
\label{sec:disc_rheology}

The fits constrain the composite parameter $a = EH/[2\eta(1-\nu^2)]$, not
the viscosity alone.  For fixed $E$, the best-fit $a$ maps onto an iso-$a$
band in the $(H,\eta)$ plane
(Figures~\ref{fig:rheology_1sigma_krafla} and~\ref{fig:rheology_1sigma_asal}),
and an independent estimate of the seismogenic thickness,
$H \approx 5$--$15$~km in Iceland and the Afar
\citep{stevens2021tectonic,plasman2017lithospheric,weinstein2017fault},
turns it into a viscosity.  For the canonical $H = 10$~km, $E = 40$~GPa and
$\nu = 0.25$, the EV RSS gives $\eta \approx 4.3\times10^{18}$~Pa~s at
Krafla and $\eta \approx 2.1\times10^{19}$~Pa~s at Asal.  The 1$\sigma$
ranges from the $\chi^2$ landscapes alone are narrow,
$[3.7,\,5.1]\times10^{18}$ and $[2.05,\,2.21]\times10^{19}$~Pa~s.  Adding
an uncertainty on $E$ over $[20,\,80]$~GPa in quadrature widens them to
$[2.1,\,8.7]\times10^{18}$ and $[1.1,\,4.3]\times10^{19}$~Pa~s: the
elastic modulus, not the fit, dominates the error budget.  The Krafla value
lies below the bound $\eta < 10^{19}$~Pa~s of
\citet{sigmundsson1991postglacial} and within the post-glacial rebound
range for Iceland \citep{sigmundsson2010geodetic}, consistent with the hot
sub-lithospheric mantle of the Iceland plume \citep{holtzman2016questions}.
The Asal value is five times higher, and higher than estimates for the
Dabbahu segment of the same rift system
($\eta \approx 5\times10^{18}$~Pa~s; \citealt{nooner2009post};
$\eta \sim 10^{18}$~Pa~s for $H = 21$~km; Hamling et al., 2014
[CITE NEEDED]).

The deformation half-width explains why the profiles constrain $b$ at all
(Figure~\ref{fig:css_width}).  With the derived recurrence intervals,
$aT \approx 149$~km at Krafla and $\approx16$~km at Asal, comparable to the
extent of each network on either side of the rift.  Each network therefore
samples the transition from the near-rift transient to the far-field
plateau $\pm V/2$, rather than the transient alone or the plateau alone, and
through $T = b/V$ this transition fixes the opening.

The fits also expose what the model leaves out.  At Krafla, the derived
$T = 89$--$95$~yr is about a third of the $\sim250$~yr between the
1724--1729 and 1975--1984 crises.  At Asal, the fitted opening
$b = 0.67$~m is about a third of the $\approx1.9$~m measured across the
rift after the 1978 event \citep{ruegg1979geodetic}.  The predicted early
lengthening rate exceeds the trilateration rate by 27 per cent
\citep{ruegg1987deformation}, and the profile retains an SW--NE asymmetry
that a symmetric model cannot produce ($\chi^2/N \approx 21$).  Each misfit
points to an ingredient absent from the model: a single, infinitely thin
dike in one elastic layer, symmetric about the axis, with no shallow
magmatic source and strictly periodic recurrence.
\citet{smittarello2016transient} argue that magma supply rather than
visco-elastic relaxation controls the post-diking velocity field at Asal,
and our single-parameter description cannot separate the two without an
independent constraint on the magmatic budget.  The fitted $(a,b)$ are
therefore effective values for a symmetric, periodic rift, and the
viscosities derived from them inherit this simplification.

\begin{figure}
  \centering
  \includegraphics[width=\columnwidth]{fig/fig_rheology_1sigma.png}
  \caption{Rheological trade-off for Krafla, for the EV RSS best-fit
    $a = 4.99\times10^{-5}$~m~s$^{-1}$ (Section~\ref{sec:krafla}).  The
    panel shows the iso-$a$ contour in the $(H,\,\eta)$ plane at the
    best-fit value.  Dark shading propagates the 1$\sigma$ uncertainty on $a$
    from the joint $\chi^2$ landscape at fixed $E = 40$~GPa; light shading additionally
    propagates $E$ over $[20,\,80]$~GPa (combined 1$\sigma$ envelope).  The
    $H = 10$~km dashed line marks the reference elastic thickness, and the
    filled circle with error bar the best fit at that thickness with its
    1$\sigma$ range at $E=40$~GPa.  Published estimates: dashed boxes,
    the channel-model viscosities of \citet{foulger1992post} and
    \citet{heki1993plate} for an elastic layer of 8--30~km (viscosities of
    a thin channel, not of a half-space, so not directly comparable); open
    square, \citet{hofton1996post}; open circle with error bar,
    \citet{wright2012geophysical}, drawn at $H=11$~km for legibility; open
    triangle, the upper bound $\eta\le10^{19}$~Pa~s of
    \citet{sigmundsson1991postglacial}.}
  \label{fig:rheology_1sigma_krafla}
\end{figure}

\begin{figure}
  \centering
  \includegraphics[width=\columnwidth]{fig/fig_rheology_1sigma_asal.png}
  \caption{Rheological trade-off for Asal, for the EV RSS best-fit
    $a = 1.00\times10^{-5}$~m~s$^{-1}$ (Section~\ref{sec:asal}).  The panel
    shows the iso-$a$ contour in the $(H,\,\eta)$ plane at the best-fit
    value.  Dark shading propagates the 1$\sigma$ uncertainty on $a$ from the
    joint $\chi^2$ landscape at fixed $E = 40$~GPa; light shading additionally
    propagates $E$ over $[20,\,80]$~GPa (combined 1$\sigma$ envelope).  The
    $H = 10$~km dashed line marks the reference elastic thickness; markers show
    published viscosity estimates: filled diamond, \citealt{nooner2009post}, and
    star with error bar, \citealt{wright2012geophysical}, both for the Dabbahu
    segment of the same rift system; orange box, \citealt{cattin2005numerical} — near-axis effective
    viscosity $\eta = 0.5$--$1\times10^{18}$~Pa~s from a 2D thermo-mechanical
    FEM model of the first 6--8~yr post-rifting at Asal, with the seismogenic
    layer thickness $H = 3$--$12$~km from the same study.  The Cattin et al.\
    box lies well below the iso-$a$ curve, consistent with the local
    near-axis viscosity being lower than the laterally averaged bulk estimate
    from the 25-yr geodetic record.}
  \label{fig:rheology_1sigma_asal}
\end{figure}

\begin{figure}
  \centering
  \includegraphics[width=\columnwidth]{fig/fig_css_width.png}
  \caption{RSS deformation half-width $aT$ as a function of recurrence
    interval $T$ for Krafla (red, $a = 4.99\times10^{-5}$~m~s$^{-1}$) and
    Asal (blue, $a = 1.00\times10^{-5}$~m~s$^{-1}$), the EV RSS best fits.
    Vertical dashed lines mark the recurrence intervals derived from
    $T=b/V$ (95~yr and 52~yr); stars give the implied $aT$.  Horizontal
    dash-dot lines mark the total span of each geodetic network used in
    this study: 230~km for Krafla \citep[distance between outermost
    stations;][]{foulger1992post} and 17~km for Asal \citep[GPS profile
    extent;][]{vigny2007twenty}.  The half-width $aT$ is comparable to the
    extent of each network on either side of the rift.}
  \label{fig:css_width}
\end{figure}

% ── §4.5 ─────────────────────────────────────────────────────────────────────
\subsection{Irregular recurrence}
\label{sec:disc_stochastic}

The RSS assumes that intrusions recur at a fixed interval $T$, with a fixed
opening $b$, always at the same place, and none of this holds in nature.
Magmatic rifting is episodic; the $\sim250$~yr gap between the Krafla
crises is a single estimate from a sparse historical record, dike openings
vary within and between sequences \citep{thordarson2007volcanism}, and
successive dikes need not open on the same plane.  We therefore tested how
closely the RSS describes a rift whose timing, opening and dike position
vary from cycle to cycle, keeping $T = b/V$ on average in every case.  Each
realisation sums the EV step velocities (\ref{eq:stepv}) of 300 past
intrusions.  We considered four cases (Figure~\ref{fig:stochastic_css}):
(a) openings $b_j$ drawn at random with $T$ fixed, a contrast in which the
opening does not match the interval; (b) intervals $T_j$ drawn at random
with each opening coupled to the preceding interval, $b_j = V T_j$;
(c) the position $x_j$ of each intrusion drawn from a Normal distribution
of standard deviation $\sigma_x = 0.1\,aT$, with periodic timing and
opening; and (d) cases (b) and (c) combined.  $T_j$ and $b_j$ follow Normal
distributions with coefficient of variation $c_v = 0.30$, truncated at
5 per cent of the mean.  Each case comprises 500 realisations at phase
$\tau/T = 0.5$.  A uniform distribution of $x_j$ with the same standard
deviation gives bands within 10 per cent of the Normal case.

In all four cases the ensemble median follows the deterministic RSS to
within 3 per cent of the far-field velocity $V/2$.  For scatter in $b_j$ or
$x_j$ alone, linearity fixes the ensemble mean: it equals the RSS for
$b_j$, and the RSS convolved with the distribution of $x_j$ for $x_j$.
Scatter in $T_j$ enters non-linearly through the ages of past intrusions,
but the resulting bias is small at $c_v = 0.30$.  The RSS therefore
describes the typical velocity field of an irregular rift, and the bands
measure how far a single realisation departs from it.

The three sources of scatter act at different distances from the axis,
because each station integrates a different set of past intrusions.  Near
the axis, the contribution of an intrusion decays with its age $t$ as
$1/t^2$ (Equation~\ref{eq:stepv}), so the most recent ones dominate;
farther out, a station also feels every intrusion younger than about
$|x|/a$.  In units of $V/2$, the 10--90 per cent band of the contrast case
(a) reaches 0.46 at $\xi = 0.3$ and still 0.20 at $\xi = 5$: when openings
do not match the intervals, the plateau of a single realisation reflects
the opening rate realised over the last few cycles.  Coupling the openings
to the intervals (b) removes most of that far-field spread (0.31 at
$\xi = 0.3$, 0.06 at $\xi = 2$, 0.01 at $\xi = 5$), because the opening
delivered over any time window then equals $V$ times that window.  Scatter
in dike position (c) acts only near the axis, where the band reaches 0.62
of $V/2$, and decays to 0.24 at $\xi = 0.3$ and below 0.01 at $\xi = 5$:
far from the axis, a shift of the source by $\sigma_x$ barely changes the
distance to the station.  Combined, coupled timing and dike position (d)
give a band of 0.6 at the axis that narrows to 0.18 at $\xi = 1$ and
0.02 at $\xi = 5$.

These results bear directly on the fits of Section~\ref{sec:results}.  Both
networks lie within one half-width of the axis, $|x| \lesssim 0.9\,aT$ at
Krafla and $|x| \lesssim 0.6\,aT$ at Asal, where the bands of the coupled
and dike-position cases are widest.  The effect of dike position scales
with $\sigma_x/(aT)$, so it matters more at Asal than at Krafla: a dike
zone a few kilometres wide corresponds to $\sigma_x/(aT)$ of order 0.1 at
Asal ($aT \approx 16$~km) but of order 0.01 at Krafla ($aT \approx 149$~km).
The fitted $(a,b)$ therefore carry an uncertainty from cycle-to-cycle
variability that the $\chi^2$ landscapes do not include.  The bands of
Figure~\ref{fig:stochastic_css} are computed at one phase and one level of
scatter, so they indicate the size of this effect rather than quantify it
for either site.

\begin{figure}
  \centering
  \includegraphics[width=\columnwidth]{fig/fig_stochastic_css.png}
  \caption{EV RSS velocity profile under stochastic recurrence
    ($a = T = b = 1$; phase $\tau/T = 0.5$; 300 past intrusions; 500
    realisations; $T = b/V$ on average in every case).  Velocity is
    normalised by $b/T$ and distance by $aT$.  Shaded bands show the
    10th--90th percentile range; coloured lines are ensemble medians; the
    black dashed line is the deterministic RSS (\ref{eq:vcss}); vertical
    dashed lines mark $\xi = \pm 1$.  (a)~Contrast: openings $b_j$ random
    ($c_v = 0.30$) with $T$ fixed.  (b)~Intervals $T_j$ random ($c_v = 0.30$)
    with each opening coupled to the preceding interval, $b_j = V T_j$.
    (c)~Dike position $x_j$ random, Normal with $\sigma_x = 0.1\,aT$, periodic
    timing and opening.  (d)~Cases (b) and (c) combined.  Insets in (c) and
    (d) enlarge $|\xi| \le 0.5$.  In all four cases the ensemble median stays
    within 3 per cent of $V/2$ of the deterministic RSS.}
  \label{fig:stochastic_css}
\end{figure}

% ── §4.6 ─────────────────────────────────────────────────────────────────────
\subsection{Perspectives}
\label{sec:disc_perspectives}

The RSS velocity depends on three combinations of the model parameters:
the half-width $aT$, the far-field velocity $b/(2T) = V/2$, and the phase
$\tau/T$ (Equation~\ref{eq:vcss}).  Our fits derive the recurrence
interval from the fitted opening and an independent rift opening rate,
$T = b/V$ (Section~\ref{sec:results}).  Records that span a larger
fraction of a cycle, or snapshots repeated over decades, would sample the
evolution of the profile with phase and so test this derived $T$ against
geodesy.  Propagating the stochastic spread of
Section~\ref{sec:disc_stochastic} into the inferred parameters, with
ensembles conditioned on each site's phase and intrusion history, would
complete the uncertainty budget.

Real rifts are segmented: dikes span finite along-strike lengths, and
segments differ in recurrence history, plate thickness and heat flow.
Because the only length scale of the model is the local product $aT$, each
segment generates its own deformation halo, and to the extent that the
segments respond linearly, the surface velocity is their superposition.
The 1D model cannot capture along-strike dike propagation or the interaction
between segments.  Extending the RSS to the full 2D problem requires
relaxing the plane-strain approximation while retaining the non-local
Stokes coupling at the base of the plate.

% ====================================================================
\section{Conclusions}
\label{sec:conclusions}

Inter-rifting deformation at magmatic rifts can be described as the
superposed relaxation of a long history of discrete dike intrusions in an
elastic plate over a viscous half-space.  In the limit of many intrusions,
this superposition reaches a Recurrent Steady State with a closed form in
the Digamma and log-Gamma functions.  The RSS conserves the opening of
each intrusion over one cycle, so far-field plate divergence follows from
the model itself rather than from a locking depth, a prescribed domain
width or a boundary condition, and its only length scale is the
deformation half-width $aT$.  A steady velocity far from the rift is
therefore the memory of intrusions that occurred many cycles earlier, not
evidence of continuous creep at depth.  This memory also sets the validity
of the RSS: a station at distance $|x|$ is described by it only if the rift
has hosted more than about $13\,|x|/(aT)$ intrusions.

The coupling of the plate to the half-space through the Stokes kernel
turns post-diking relaxation into a fractional diffusion governed by the
half-Laplacian.  Deformation then spreads from the rift at the constant
speed $a$ with a $1/x$ tail, instead of the decelerating front and Gaussian
tail of Fickian diffusion.  The Fickian model of \citet{foulger1992post} is
the same relaxation evaluated at a single diffusion time, whereas the EV
solution averages Fickian profiles over a heavy-tailed distribution of
diffusion times.  At Krafla, this difference is visible in the data: under
the same conditions, the EV RSS fits the two post-rifting profiles with
$\chi^2_r = 1.61$, against 4.90 for the diffusion model.

A Maxwell half-space changes the response only within the Maxwell length
$\ell_M$ of the axis and during the first few relaxation times $\tau_R$.
With the plate and the half-space sharing one elastic modulus,
$\ell_M = H/(1-\nu) \approx 1.3\,H$, so a geodetic record can distinguish
the two rheologies only if it starts early in units of $\tau_R$ and has
stations close to the axis.  The Krafla record, which starts
$3.4\,\tau_R$ after the onset, favours the viscous half-space only
modestly ($\Delta\mathrm{AIC} = +7.1$); the Asal record, which starts at
$0.2\,\tau_R$, favours it decisively ($\Delta\mathrm{AIC} = +70.3$).  The
choice of rheology also changes the inferred viscosity, by a factor of up
to seven at Asal.

Tying the recurrence interval to an independent rift opening rate,
$T = b/V$, leaves two free parameters, $a$ and $b$ (plus a reference-frame
offset at Krafla), which the Krafla and Asal records both constrain.  For $H = 10$~km and $E = 40$~GPa, the EV RSS
implies $\eta \approx 4.3\times10^{18}$~Pa~s at Krafla and
$\eta \approx 2.1\times10^{19}$~Pa~s at Asal, with an error budget
dominated by the elastic modulus rather than by the fit.  The remaining
misfits, a recurrence interval shorter than the Krafla eruptive history and
an opening and a profile asymmetry at Asal that the model does not
reproduce, point to the simplifications of the model: a single,
infinitely thin dike in one elastic layer, symmetric about the axis, with
no shallow magmatic source and strictly periodic recurrence.

Irregular recurrence spreads individual velocity profiles around the RSS
with little bias: with 30 per cent scatter in recurrence interval and
opening, and dike positions scattered over $0.1\,aT$, the ensemble median
stays within 3 per cent of the deterministic RSS.  The spread is largest
within about one half-width of the axis, where the velocity reflects the
last few intrusions.  Scatter in dike position acts only within a few
$\sigma_x$ of the axis, and the far-field spread nearly vanishes when each
opening scales with the preceding interval, as $T = b/V$ requires.  Extending the framework
to segmented rifts and to the full along-strike 2D geometry, while
retaining the non-local coupling of the plate to the half-space, is the
next step toward a system-scale model of strain accumulation at magmatic
plate boundaries.

% ====================================================================
\section*{Acknowledgements}
[PLACEHOLDER — funding agencies, field/computing support, reviewer thanks.]

% ====================================================================
\section*{Data Availability}
The geodetic data used in this study are from published sources cited in the
text. The analytical model code and inversion scripts are available at
[PLACEHOLDER — repository DOI or URL upon acceptance].

\bibliographystyle{gji}
\bibliography{refs}

\appendix

\section{Fourier transform of the step initial condition}
\label{app:step}

The Fourier transform of the initial condition (\ref{eq:IC}) is
$\hat{u}(k,0) = b\int_{-\infty}^{0}e^{-ikx}\,dx$, which for real $k$ is
treated in the distributional sense by adding a convergence regulator
$\varepsilon > 0$:
\begin{equation*}
\int_{-\infty}^{0}e^{-ikx}\,dx
= \lim_{\varepsilon\to 0}\int_{-\infty}^{0}e^{(-ik+\varepsilon)x}\,dx
= \lim_{\varepsilon\to 0}\frac{1}{\varepsilon - ik}.
\end{equation*}
By the Sokhotski--Plemelj theorem \citep{ablowitz2003complex},
\begin{equation}
\lim_{\varepsilon\to 0^+}\frac{1}{x \pm i\varepsilon} = \mathcal{P}\left(\frac{1}{x}\right) \mp i\pi\delta(x),
\label{eq:sokhotski}
\end{equation}
where $\mathcal{P}$ denotes the Cauchy principal value. Multiplying numerator
and denominator by $i$ recasts the regulator as $+i\varepsilon$,
\begin{equation*}
\frac{1}{\varepsilon - ik} = \frac{i}{i\varepsilon + k} = \frac{i}{k + i\varepsilon},
\end{equation*}
so that (\ref{eq:sokhotski}) applies directly with $x = k$ and the upper sign:
\begin{align}
\lim_{\varepsilon\to 0^+}\frac{1}{\varepsilon - ik}
&= \lim_{\varepsilon\to 0^+}\frac{i}{k + i\varepsilon} \notag\\
&= i\left[\mathcal{P}\left(\frac{1}{k}\right) - i\pi\delta(k)\right] \notag\\
&= \pi\delta(k) + \frac{i}{k}, \label{eq:sokhotski_applied}
\end{align}
where the principal-value symbol on $1/k$ is dropped, as is standard, since
it is understood implicitly whenever $1/k$ appears next to a $\delta(k)$
term. This gives (\ref{eq:IC_FT}).

The inverse transform of the evolved spectrum (\ref{eq:evolved}) splits into
two contributions, $u(x,t) = I_1 + I_2$. The delta-function term gives a
constant background $I_1 = b/2$; for $I_2$, the parity of $\sin(kx)/k$
yields
\begin{equation}
I_2 = -\frac{b}{\pi}\int_0^{\infty}\frac{\sin(kx)}{k}\,e^{-akt}\,dk
    = -\frac{b}{\pi}\arctan\!\left(\frac{x}{at}\right),
\label{eq:I2}
\end{equation}
a standard Dirichlet-type integral. With the complementary arctangent
identity $\arctan(z)+\arctan(1/z)=\pi/2$, the displacement counted from the
instant after the intrusion is the step solution (\ref{eq:step}).

\section{Digamma representation of the RSS velocity}
\label{app:digamma}

Each term of the RSS series (\ref{eq:css_sum}) splits into partial
fractions,
\begin{equation*}
\frac{x}{x^2+a^2(\tau+jT)^2}
= \frac{1}{2i}\!\left[
    \frac{1}{a(\tau+jT)-ix} - \frac{1}{a(\tau+jT)+ix}
  \right].
\end{equation*}
Factoring $aT$ from each denominator and substituting into
(\ref{eq:css_sum}), the sum takes the form of a difference of two Digamma
series. With the representation $\psi(z) = -\gamma +
\sum_{j=0}^\infty[1/(j+1) - 1/(j+z)]$ and $z=\tau/T\pm ix/(aT)$, the
Euler--Mascheroni constant $\gamma$ and the $1/(j+1)$ terms cancel when the
two conjugate series are differenced. With $Z = \tau/T + ix/(aT)$ and the
identity $\psi(Z)-\psi(\bar{Z}) = 2i\,\IM\{\psi(Z)\}$, the series reduces to
(\ref{eq:vcss}).

\section{Correspondence principle for the Maxwell half-space}
\label{app:correspondence}

For deviatoric stress $\tau_{ij}$ and strain $\varepsilon$, the 1D
constitutive relation of a Maxwell body is
\begin{equation}
\dot{\varepsilon} = \frac{\dot{\tau}}{2\mu_m} + \frac{\tau}{2\eta}.
\label{eq:maxwell_constitutive}
\end{equation}
Applying the Laplace transform $\tilde{f}(s)=\mathcal{L}\{f(t)\}$ to
(\ref{eq:maxwell_constitutive}) with zero initial conditions gives
\begin{equation}
s\tilde{\varepsilon} = \tilde{\tau}\!\left(\frac{s}{2\mu_m}+\frac{1}{2\eta}\right).
\label{eq:laplace_ve}
\end{equation}
Solving for stress in terms of the strain rate $s\tilde{\varepsilon}$ gives
$\tilde{\tau} = 2\tilde{\eta}(s)\,(s\tilde{\varepsilon})$, with the effective
viscosity $\tilde{\eta}(s)$ of (\ref{eq:eta_eff}). Laplace-transforming the
Stokes kernel (\ref{eq:stokes}) and replacing $\eta$ by $\tilde{\eta}(s)$
gives
\[
\tilde{\sigma}_{xz}^{\text{BOT}}(x,s) =
  \int_{-\infty}^{+\infty}\frac{2\tilde{\eta}(s)}{\pi(x-x')^2}
  \left[s\tilde{u}(x',s)\right]dx',
\]
or, with (\ref{eq:eta_eff}),
\begin{equation}
(1+\tau_R s)\,\tilde{\sigma}_{xz}^{\text{BOT}}(x,s) =
  \int_{-\infty}^{+\infty}\frac{2\eta}{\pi(x-x')^2}
  \left[s\tilde{u}(x',s)\right]dx'.
\label{eq:ve_kernel_laplace}
\end{equation}
Since multiplication by $s$ in the Laplace domain corresponds to
$\partial/\partial t$, (\ref{eq:ve_kernel_laplace}) inverts to the
visco-elastic Stokes kernel (\ref{eq:ve_kernel}).

\section{Numerical evaluation of the EVE integrals}
\label{app:numerics}

The EVE step integrals (\ref{eq:maxwell_step}) and (\ref{eq:maxwell_vel})
are evaluated with adaptive Gauss--Kronrod quadrature
\citep{laurie1997calculation}, with integration breakpoints at the
successive zeros of $\sin(kx)$, $k = n\pi/x$ ($n = 1,2,\ldots$), which
partition the oscillatory integrand into well-behaved sub-intervals. The
upper truncation wavenumber is set adaptively so that the exponential factor
is negligible at the cutoff.

The EVE RSS integrals (\ref{eq:maxwell_css_v}) and (\ref{eq:maxwell_css_u})
contain factors that tend to finite, non-zero limits as $k\to\infty$, so
their integrands decay slowly. For the fits of Section~\ref{sec:results},
the large-$k$ limit of each integrand is subtracted and integrated in closed
form with the Dirichlet integral, which also gives the axis jump
(\ref{eq:maxwell_rss_jump}). The remainder decays as $1/k^2$ on the scale
$k \sim 1/(a\tau_R)$ and is integrated with the same segmented quadrature
up to the fixed cutoff $10^4/(a\tau_R)$.


\label{lastpage}
\end{document}


